\documentclass[11pt]{article}

\usepackage[utf8]{inputenc}
\usepackage[T1]{fontenc}
\usepackage{lmodern}
\usepackage{booktabs}
\usepackage{enumitem}
\usepackage{amsmath, amssymb, amsthm}
\usepackage[margin=1in]{geometry}
\usepackage{microtype}
\PassOptionsToPackage{hyphens}{url}
\usepackage{xcolor}
\usepackage{hyperref}
\hypersetup{
  colorlinks=true, linkcolor=blue!50!black, urlcolor=blue!50!black, citecolor=blue!50!black,
  pdftitle={The growth series of the generic right-triangle reflection group is not D-finite},
  pdfauthor={Vico Bonfioli},
}

\theoremstyle{plain}
\newtheorem{theorem}{Theorem}
\newtheorem{lemma}[theorem]{Lemma}
\newtheorem{corollary}[theorem]{Corollary}
\newtheorem{proposition}[theorem]{Proposition}
\theoremstyle{definition}
\newtheorem{definition}[theorem]{Definition}
\newtheorem{hypothesis}[theorem]{Hypothesis}
\theoremstyle{remark}
\newtheorem{remark}[theorem]{Remark}

\newcommand{\Z}{\mathbb{Z}}
\newcommand{\Q}{\mathbb{Q}}
\newcommand{\R}{\mathbb{R}}
\newcommand{\C}{\mathbb{C}}
\newcommand{\F}{\mathbb{F}}
\renewcommand{\Re}{\operatorname{Re}}
\theoremstyle{plain}\newtheorem*{maintheorem}{Main Theorem}

\title{The growth series of the generic right-triangle reflection group\\ is not $D$-finite}
\author{Vico Bonfioli\thanks{Independent researcher. \texttt{vicobonfioli@gmail.com}.
Code, data, and the Lean~4 formalisation:
\url{https://github.com/ElVec1o/pythagorean-reflection-sequence}.}}
\date{2026}

\begin{document}
\maketitle

\begin{abstract}
\noindent Let $W_{\mathrm{univ}}=\langle R_0,R_1,R_2\rangle$ be the group generated by the
reflections in the sides of a right triangle with generic angle parameter, to which every right
triangle group agrees up to a finite, effectively computable depth \cite{Bonfioli2026a}. Let
$U(x)=\sum_nu_nx^n$ be its growth series with respect to the three side reflections (the integer
sequence \textup{OEIS~A396406} \cite{OEIS}), and $V$ its relaxed companion. We prove that $U$ and $V$
continue meromorphically to the disc $|x|<1$ and have infinitely many poles there, accumulating at
the unit circle. Consequently neither is algebraic over $\Q(x)$ or $D$-finite, so $(u_n)$ and
$(v_n)$ are not P-recursive; neither satisfies a linear $q$-difference equation with $Q$ not a root
of unity or a Mahler equation; and, by the P\'olya--Bertrandias theorem, $|x|=1$ is a natural
boundary of both. The poles lie over the zeros of the Hahn--Exton $q$-Bessel function
$J^{(3)}_{-1/2}(Z(q);q^2)$, $Z=\sqrt{2(1-q)/q}$, $q=x^2$, which up to a positive factor is the
Fredholm determinant of the travel transfer operator; that there are infinitely many of them
follows from a spectral argument with no asymptotic estimate. They survive in $U$ and $V$ because
the leading Laurent coefficient at each of them is the square of a junction pairing, which has an
exact closed form in one bulk quantity and is positive at every travel pole. For the travel
resolvent $\Sigma_0/(1-\Sigma_1)$ all of this is unconditional. For $U$ and $V$ it rests on the
transfer model, a theorem whose assembly step is a hand proof, and on the junction-pairing
theorem, whose positivity uses an analytic gate estimate proved partly by verified interval
arithmetic. The long proofs are given in a companion paper.
\end{abstract}

\section{Introduction}\label{sec:intro}

\paragraph{Growth series of groups.} For a group $G$ with a finite generating set $S$, the growth
series $\sum_{n\ge0}a_nx^n$ counts the elements of $G$ by their word length in $S$. In the
classical families it is rational: hyperbolic groups have rational growth series with respect to
every finite generating set (Cannon \cite{Cannon1984}), and Coxeter systems are governed by
Steinberg's rational formula \cite[\S5.12]{Humphreys1990}. For wreath products, including
lamplighter-type groups, Parry computed the growth series for natural generating sets and found
them \emph{algebraic} \cite{Parry1992}. Provably transcendental growth series are rarer. Stoll
\cite{Stoll1996} showed that the higher Heisenberg group $H_5$ has a transcendental growth series
with respect to its standard generating set, and a rational one with respect to a specially chosen
generating set; the three-dimensional Heisenberg group, by contrast, has rational growth for every
finite generating set (Duchin--Shapiro \cite{DuchinShapiro2019}). Two general mechanisms give growth
series that are not even $D$-finite. For groups of polynomial growth, and for groups of
intermediate growth such as Grigorchuk's \cite{Grigorchuk1984}, the growth series has radius of
convergence $1$ and integer coefficients, so by P\'olya--Carlson \cite{Carlson1921} it is rational
or has $|x|=1$ as a natural boundary; in the second case, which includes Stoll's series for $H_5$
and every group of intermediate growth, it is not $D$-finite. And a finitely presented group with
unsolvable word problem has a non-computable growth function, while a $D$-finite series with
rational coefficients has computable coefficients, so its growth series is not $D$-finite either.

\paragraph{The group.} A right triangle $T$ has three sides; the reflections across them generate a
group $W_T=\langle R_0,R_1,R_2\rangle\le\mathrm{Aff}(\R^2)$. For a right triangle whose angle
parameter is generic, $W_T$ is the \emph{universal} group $W_{\mathrm{univ}}$ of
\cite{Bonfioli2026a}, and every other right triangle group agrees with $W_{\mathrm{univ}}$ in the
ball of some finite radius and deviates beyond it; the radius is effective
\cite[Theorem ``Effective universality'', \textup{thm:effective-universality}]{Bonfioli2026a}
\textup{(}for the $(1,2)$ triangle the first deviation is at depth $33$,
\cite[Theorem ``Sharp uniform universality at depth $32$'', \textup{thm:universality-sharp}]{Bonfioli2026a}\textup{)}.
Irrationality of the angles alone \textup{(}Niven \cite{Niven1956}\textup{)} does not give
genericity: a triangle in $\Q^2$ with unequal legs has irrational angles, yet it is not universal.
Throughout, $U(x)=\sum_nu_nx^n$ is the growth series of $W_{\mathrm{univ}}$ with respect to
$\{R_0,R_1,R_2\}$. The sequence $u_n=1,3,5,8,13,21,34,55,89,\dots$ is \textup{OEIS~A396406}; it
coincides with the Fibonacci numbers $F_{n+3}$ for $1\le n\le9$ and deviates at $n=10$.
$W_{\mathrm{univ}}$ is not a Coxeter group \cite{IsolaPiergallini}, and it is virtually the
lamplighter $\Z\wr\Z$ \cite{Bonfioli2026a}. Its generating set is not constructed but imposed by
Euclidean geometry. For a group generated by reflections, Steinberg's formula would make the
growth series rational in the Coxeter case, and Parry's theorem makes the natural wreath-product
series algebraic; the result below places this geometric system outside both classes.

The word length $\ell_T(g)$ is realised by strand walks in a finite transition system; permitting
free isolated cycles yields a smaller \emph{relaxed} length $\ell_R(g)$, with
$\ell_T=\ell_R+2c(g)$, $c(g)\in\Z_{\ge0}$ a cycle count. The relaxed counting function is the
companion $V$; one has $v_n\ge u_n$, with equality for $n\le7$ and strict inequality at $n=8$.

\paragraph{The main theorem.}

\begin{maintheorem}
\begin{enumerate}[label=\textup{(\roman*)},leftmargin=2.2em,itemsep=0.1em,topsep=0.2em]
\item $U$ and $V$ continue to single-valued meromorphic functions on the disc $|x|<1$, and each has
infinitely many poles there, accumulating at the unit circle
\textup{(}Lemma~\ref{lem:mero}, Theorems~\ref{thm:V} and~\ref{thm:U}\textup{)}.
\item Consequently $U$ and $V$ are transcendental over $\Q(x)$ and not $D$-finite, so $(u_n)$ and
$(v_n)$ are not P-recursive \textup{(}Theorem~\ref{thm:nonDfinite}\textup{)}; neither satisfies a
linear $q$-difference equation $\sum_ia_i(x)f(Q^ix)=b(x)$ with $Q$ not a root of unity, nor a Mahler
equation \textup{(}Proposition~\ref{prop:noqdiff}\textup{)}; and $|x|=1$ is a natural boundary of
both \textup{(}Theorem~\ref{thm:natboundary}\textup{)}.
\item The travel resolvent $G^T_0=\Sigma_0/(1-\Sigma_1)\in\Z[[q]]$ has all of these properties in
the variable $q$, with natural boundary $|q|=1$ \textup{(}Theorems~\ref{thm:travelres},
\ref{thm:nonDfinite}, \ref{thm:natboundary}, Proposition~\ref{prop:noqdiff}\textup{)}.
\item The exponential growth rate of $(u_n)$ and of $(v_n)$ is $\beta_2=1.4916177871\ldots$, where
$\beta_2^{-2}$ is the smallest zero in $(0,1)$ of $q\mapsto J^{(3)}_{-1/2}(Z(q);q^2)$
\textup{(}Corollary~\ref{cor:beta}\textup{)}.
\end{enumerate}
\textup{(}Standing: \textup{(iii)} is unconditional. \textup{(i)}, \textup{(ii)} and \textup{(iv)}
rest on Theorems~\ref{thm:model} and~\ref{thm:RJ}, as stated there.\textup{)}
\end{maintheorem}

\paragraph{The mechanism.} The proof does not go through the coefficients. It locates infinitely
many poles of $U$ and $V$ inside the disc $|x|<1$ of meromorphy, beyond the radius of convergence
$1/\beta_2=0.6704\ldots$, and then applies three classical facts: an algebraic function has
finitely many poles in the disc; a pole of a solution of a linear differential equation whose
coefficients are holomorphic near the closed disc is a zero of its leading coefficient
(Lemma~\ref{lem:odeposes}); and an integer power series that is meromorphic on a simply connected
domain of conformal radius larger than $1$ is rational (P\'olya--Bertrandias,
Theorem~\ref{thm:PB}). The poles are produced in three steps.

\emph{(1) The transfer model.} Grading by relaxed length and marking cycles gives a bivariate
series $\mathcal W(x,y)$ with $V=\mathcal W(x,1)$ and $U=\mathcal W(x,x^2)$. Theorem~\ref{thm:model}
writes $\mathcal W$ as a sum of three sectors. Two of them are pairings against the resolvent
$(I-\mathcal T)^{-1}$ of the \emph{travel operator} $\mathcal T=D_eK$, with $q=x^2$,
$K=[q^{\max(s,s')}]_{s,s'\ge0}$ and $D_e=\operatorname{diag}(2q^{1+s})$; the bulk enters only
through explicit vectors built from the resolvent of the same semiseparable kernel on a shifted
index set, with a rank-one gap term, and the third sector contains no travel resolvent.
Lemma~\ref{lem:mero} continues $\mathcal W$ meromorphically to $|x|<1$.

\emph{(2) The travel poles.} The poles of $(I-\mathcal T)^{-1}$ are the zeros of its Fredholm
determinant, and Theorem~\ref{thm:fredholm} computes that determinant in closed form:
\[
   \det(I-\lambda\mathcal T)=\sum_{k\ge0}\frac{(-2\lambda(1-q))^kq^{k^2}}{(q;q)_{2k}}
   =\cos(\sqrt\lambda\,Z;q^2)=c_q\,(\sqrt\lambda\,Z)^{1/2}J^{(3)}_{-1/2}(\sqrt\lambda\,Z;q^2),
\]
the $q$-cosine of Koornwinder and Swarttouw \cite{KS1992} and the Hahn--Exton $q$-Bessel
function \cite{Hahn1953,Exton1978,KoelinkSwarttouw1994}. So the travel poles are the $q\in(0,1)$
with $J^{(3)}_{-1/2}(Z(q);q^2)=0$ (Corollary~\ref{cor:qbessel}). A symmetrisation of $\mathcal T$ is
a positive trace-class operator with simple eigenvalues $1/\lambda_k(q)$; the branches satisfy
$\lambda_1(q)>(1-q^2)/(2q)$ and $\lambda_k(q)/(1-q)\to(2k-1)^2\pi^2/8$, so each crosses the level
$1$, and distinct branches cross at distinct points: there are infinitely many travel poles
(Proposition~\ref{prop:infpolesalt}), locally finite in $(0,1)$ (Lemma~\ref{lem:polefinite}), hence
accumulating at $q=1$. This step uses no asymptotic estimate of a $q$-series.

\emph{(3) The residue.} At a travel pole $q_m$ the leading Laurent coefficient of
$\mathcal W(\cdot,y)$ is $c_\star\Pi_y(q_m)^2/(2x_m)$, where $\Pi_y=\langle\lambda^{(y)},R\rangle$
pairs the junction vector with the null vector $R$ of $I-\mathcal T(q_m)$
(Proposition~\ref{prop:shape}). The junction is not of rank one, so the numerator does not factor
through the total bulk sum; nevertheless two shift identities for $R$ collapse the pairing to an
exact closed form in the single bulk quantity $t_1$ (Theorem~\ref{thm:RJ}), for instance
\[
   (1-\mathfrak g_Vt_1)\,\Pi_1(q_m)=\frac{2q(1+x)}{1-q}\Bigl(t_1+\frac{1+x}{2x}\Bigr),
\]
and this is positive at every travel pole: analytically for $\tau_m\le5\cdot10^{-3}$ once the
\emph{gate} $|\mathfrak g_Vt_1|<1$ holds, and at the first twelve poles by an interval certificate.
For the travel resolvent the corresponding non-vanishing is an exact identity,
$\Sigma_0(q_m)S_0(q_m)=2q_m/(1-q_m)$ (Proposition~\ref{prop:numnonvanish}), a consequence of the
$q$-Pythagorean identity at a zero of the $q$-cosine.

\paragraph{Standing, and what is cited from the companion paper.} The statements about the travel
resolvent are unconditional. The statements about $U$ and $V$ rest on two theorems. The transfer
model \textup{(M)}, Theorem~\ref{thm:model}: its metric identity and the faithfulness of the model
are formalised in Lean~4, and its block decomposition \textup{(M3$'$)} is a hand proof, given in the
companion paper \cite{Bonfioli2026bB}, checked against exact enumeration to $x^{26}$ and not
formalised. The junction pairing \textup{(R-J)}, Theorem~\ref{thm:RJ}: its closed form is proved
here; its positivity is analytic for $\tau_m\le5\cdot10^{-3}$, given the gate, and rests on an
interval-arithmetic certificate at the first twelve poles; only the analytic range, which contains
all but finitely many travel poles, is needed for (i)--(ii) of the Main Theorem. The gate itself,
the amplitude estimate $(\star)$ and the lower bound on the bulk denominator at the travel poles,
is proved in \cite{Bonfioli2026bB}, with constants produced by committed certificates and
re-certified in directed-rounding interval arithmetic on the range stated there.

This paper states and proves the spectral and $q$-trigonometric steps, the local form, the closed
form of the junction pairing and every consequence above, and it states the transfer model with an
outline of its proof. The long technical proofs are in \cite{Bonfioli2026bB} and are cited by
theorem label: the block assembly (the section identity, the telescoping of the travel and bulk
sums, the reciprocity and the M\"obius factorisation of the gap-marked bulk resolvent), the site
cost of the crossing optimisation and its Lean formalisation, the gap-free theorem behind
Proposition~\ref{prop:travelinv}, the proof of \textup{(M3$'$)}, the gate with its certificates, the
explicit lower bound for $S_e$ at the travel poles, an absolute bound on an auxiliary contour
integral $T_2$ which gives a second proof of the infinitude of the travel poles and localises them
in windows of $w$ \textup{(}the window statement is used only in the Mahler part of
Proposition~\ref{prop:noqdiff}\textup{)}, the
P\'olya--Carlson dichotomy for the four building blocks, the spectral sign laws with their standing,
and the value-transcendence question for the number $\beta_2$. The archival full version
\cite{Bonfioli2026b} contains both papers; labels are shared with it.

\paragraph{Related results.} The cogrowth series is settled by general theory: $W_{\mathrm{univ}}$
is amenable, being virtually metabelian, and not virtually nilpotent, so its cogrowth series is not
$D$-finite by Bell--Mishna \cite{BellMishna2020}. In higher dimensions the picture is subtler
\cite{Bonfioli2026c}: the abstract right-angled Coxeter envelope of an $n$-orthoscheme has rational
growth, but the geometric group is an amenable proper quotient that matches this envelope only to a
finite depth. Whether $U$ and $V$ are $D$-algebraic is open.

\paragraph{Organisation.} Section~\ref{sec:setup} introduces the building blocks and the travel
poles. Section~\ref{sec:travelpoles} identifies the pole equation with the $q$-cosine and with the
Fredholm determinant of the travel operator, proves that the travel poles are infinite and
accumulate only at $q=1$, and proves that the numerators do not vanish there.
Section~\ref{sec:modelassembly} states the transfer model, continues $\mathcal W$ meromorphically
and computes the shape of its local form at a travel pole; Section~\ref{sec:RJ} evaluates the
junction pairing. Section~\ref{sec:V} proves transcendence, and Section~\ref{sec:beyond} the
statements beyond algebraicity. Section~\ref{sec:conclusion} lists open questions.

\section{The catalytic equation and its building blocks}\label{sec:setup}

Write $q=x^2=e^{-\tau}$, $w=\sqrt{2/\tau}$. The relaxed length statistic satisfies a linear
$q$-difference equation with dilation $t\mapsto q^2t$; telescoping it produces four integer power
series in $q$,
\begin{equation}\label{eq:blocks}
  \Sigma_1(q)=\sum_{i\ge1}(-1)^{i-1}\frac{w^{2i}}{(2i)!}\,\eta_i,\quad
  S_e(q)=\sum_{j\ge0}\frac{(-2(1-q))^j q^{j(j+1)}}{(q;q)_{2j}},
\end{equation}
and their companions $\Sigma_0,S_0$, where $\eta_j=e^{-B_j}$ is the Gaussian-decaying form factor
of {\cite[\S\,\textup{sec:sd}]{Bonfioli2026bB}} (equivalently $\eta_j=q^{j^2+j}[2(1-q)]^j(2j)!/[(q;q)_{2j}W^{2j}]$ with $W=we^{-\tau/2}$; the
identification of this display with $1-\Sigma_1=\cos(Z;q^2)$ of Proposition~\ref{prop:qtrig} is the
collapse used in {\cite[Lemma ``infinitude and accumulation of the travel poles'', \textup{lem:infpoles}]{Bonfioli2026bB}}). We write $\Sigma_1^{\mathrm T}=\Sigma_1$ when the travel
r\^ole is to be emphasised. The relaxed travel resolvent is $G^T_0=\Sigma_0/(1-\Sigma_1)$
({\cite[Proposition ``the travel $k$-sum telescopes'', \textup{prop:travelsum}]{Bonfioli2026bB}}), and $U$ inherits its denominator $1-\Sigma_1$
(Proposition~\ref{prop:travelinv}); how the two blocks assemble into $V$ and $U$, and how far that
assembly is proved, is {\cite[\S\,\textup{sec:assembly}]{Bonfioli2026bB}}. The \emph{travel poles}
\[
   q_1<q_2<\cdots\uparrow1,\qquad \Sigma_1^T(q_m)=1,
\]
form an infinite set accumulating at $q=1$ (Proposition~\ref{prop:infpolesalt} with
Lemma~\ref{lem:polefinite}; a second proof, through an absolute bound on an auxiliary contour
integral, is {\cite[Lemma ``infinitude and accumulation of the travel poles'', \textup{lem:infpoles}]{Bonfioli2026bB}}).

The four blocks are the $k$-recursion sums
\begin{equation}\label{eq:krec}
  \Sigma_k=\sum_{j\ge0}A_{k+2j}\prod_{i<j}C_{k+2i},\qquad
  S_k=\sum_{j\ge0}\alpha_{k+2j}\prod_{i<j}\gamma_{k+2i},
\end{equation}
\[
  A_k=\frac{2q}{1-q^{k+1}},\quad
  C_k=\frac{2q^{k+3}}{1-q^{k+2}}-\frac{2q^{k+2}}{1-q^{k+1}},\quad
  \alpha_k=\frac{2q^{k+1}}{1-q^{k+1}},\quad
  \gamma_k=\alpha_{k+1}-\alpha_k ,
\]
$\Sigma_0,\Sigma_1$ carrying the travel r\^ole and $S_0,S_1$ the bulk r\^ole; $\Sigma_1$ and $S_1$
are the denominators of \eqref{eq:blocks} and $\Sigma_0,S_0$ their numerators. Every factor is a
Lambert series with integer coefficients, and the $j$-th summand of $\Sigma_1$ starts at order
$(j+1)^2$, so each coefficient is a finite integer sum; $\Sigma_1=2q+2q^3-4q^4+6q^5-\cdots$. All
four blocks are holomorphic on $|q|<1$: by Proposition~\ref{prop:qtrig} and
Lemma~\ref{lem:qsinenum} below each is a $q$-trigonometric series whose $j$-th term carries the
Gaussian factor $q^{j^2+O(j)}$ against a denominator $(q;q)_{2j}$ or $(q;q)_{2j+1}$, bounded below
by $|(q;q)_\infty|>0$ on compact subsets of the disc \textup{(}this is the first clause of
{\cite[Theorem \textup{thm:blocks}]{Bonfioli2026bB}}, where the P\'olya--Carlson dichotomy for the four blocks is proved\textup{)}.
The coefficients of $G^T_0$ are the travel-run counts $2,6,10,26,54,114,274,\dots$, integers.

\begin{proposition}[Invariance of the travel block]\label{prop:travelinv}
The denominator $1-\Sigma_1$ is the identical, cycle-marker-free object in the generating functions
of $U$ and of $V$; hence $U$ and $V$ share the travel poles $q_m$.
\end{proposition}
\begin{proof}
A pure-travel element has no gap edge, since its span is $I_k$ and $f_j=\pm1$ there \textup{(}$I_k$ the travel interval and $f_j$ the travel indicator of
\cite[\S\,\textup{sec:sitecost}]{Bonfioli2026bB}; this is \cite[Corollary \textup{cor:localzero}]{Bonfioli2026bB}\textup{)}, so
{\cite[Theorem ``no gap edge forces $c=0$'', \textup{thm:nogap}]{Bonfioli2026bB}} gives $c=0$: some relaxed-optimal realisation has no isolated cycle. That
realisation is a single open walk, hence admissible for $\ell_T$, so $\ell_T\le\ell_R$; and
$\ell_T\ge\ell_R$ because the relaxed minimum ranges over a superset of the realisations. Therefore
$\ell_T=\ell_R$ on pure-travel elements and the two gradings of $T$ agree termwise. The travel
recursion is assembled from pure-travel runs alone, appending one $f=\pm1$ edge at a time, so it is
computed in the sector $c\equiv0$ and the travel block carries $y^0$ only.
\end{proof}

\noindent The history of this proposition, which earlier versions carried as the hypothesis
\textup{(T)}, is recorded in {\cite[\S\,\textup{sec:inputT}]{Bonfioli2026bB}}.

\section{The travel poles}\label{sec:travelpoles}

\subsection{The travel-pole equation as a $q$-trigonometric identity}\label{sec:qtrig}

The two series that govern the poles admit an exact identification with the $q$-cosine of
Koornwinder and Swarttouw \cite{KS1992}, $\cos(z;q^2):={}_1\phi_1(0;q;q^2,q^2z^2)
=\sum_{j\ge0}(-1)^j q^{j^2+j}z^{2j}/(q;q)_{2j}$.

\begin{proposition}\label{prop:qtrig}
Let $z_0=\sqrt{2(1-q)}$ and $Z=z_0/\sqrt q$. Then
\[
   1-S_1=S_e=\cos(z_0;q^2),\qquad 1-\Sigma_1=\cos(Z;q^2).
\]
In particular the travel poles $q_m$ are exactly the solutions of $\cos(Z;q^2)=0$, and $S_e$ is the
same function evaluated one half $q$-step below its own zero.
\end{proposition}

\begin{proof}
Collapse the $k$-recursion \eqref{eq:krec} on its two odd ladders, the travel one carrying
$\Sigma_1$ and the bulk one carrying $S_1$. Write $1-q^m=2e^{-m\tau/2}\sinh(m\tau/2)$ in the
$\sinh$-forms of the two series
({\cite[\S\,\textup{sec:sd}]{Bonfioli2026bB}}): the products collapse to
$1-\Sigma_1=\sum_j(-1)^j\bigl(2(1-q)\bigr)^jq^{j^2}/(q;q)_{2j}$ and
$1-S_1=\sum_j(-1)^j\bigl(2(1-q)\bigr)^jq^{j^2+j}/(q;q)_{2j}$, the second being $S_e$ as defined in
\eqref{eq:blocks}. Comparing with the series for
$\cos(z;q^2)$ term by term forces $q^2z^2=2(1-q)q$, i.e.\ $z=Z$, in the first case and
$q^2z^2=2(1-q)q^2$, i.e.\ $z=z_0$, in the second; both series converge for $|q|<1$, the factor
$q^{j^2}$ dominating. Both collapses are recomputed as identities of integer power series, to
degree $1160$ in \texttt{blocks\_growth.py} and at two truncation degrees in
\texttt{f2\_pole\_route.py}. The per-term identity is machine-checked in Lean~4
(\texttt{GramCasoratian.lean}, theorem \texttt{sigma1\_term}; axioms \texttt{propext},
\texttt{Classical.choice}, \texttt{Quot.sound}, no \texttt{sorry}).
\end{proof}

\noindent\emph{The Hahn--Exton normalisation.} Write
$J^{(3)}_\nu(x;q^2)=x^\nu\frac{(q^{2\nu+2};q^2)_\infty}{(q^2;q^2)_\infty}\,
{}_1\phi_1(0;q^{2\nu+2};q^2,q^2x^2)$ for the Hahn--Exton $q$-Bessel function, in the
normalisation of \cite[eq.~(2.4)]{KoelinkSwarttouw1994} with base $q^2$, and put
\[
   c_q:=\frac{(q^2;q^2)_\infty}{(q;q^2)_\infty}>0 .
\]
Since $(q^2;q^2)_k(q;q^2)_k=(q;q)_{2k}$ and $(q^2;q^2)_k(q^3;q^2)_k=(q;q)_{2k+1}/(1-q)$, comparing
coefficients gives, for $x>0$,
\[
   \cos(x;q^2)=c_q\,x^{1/2}J^{(3)}_{-1/2}(x;q^2),\qquad
   \sin(x;q^2)=c_q\,x^{1/2}J^{(3)}_{1/2}(x;q^2).
\]
So the travel poles are the zeros in $(0,1)$ of $q\mapsto J^{(3)}_{-1/2}(Z(q);q^2)$
(Corollary~\ref{cor:qbessel}); \cite[Proposition ``structure of the zero curve'', \textup{prop:zerostructure}(i)]{Bonfioli2026bB} records the same identity in
the variable of {\cite[Appendix \textup{app:beta2}]{Bonfioli2026bB}}.


\subsection{The pole equation as a Fredholm determinant}\label{sec:fredholmdet}

The $q$-cosine of Proposition~\ref{prop:qtrig} is the Fredholm determinant of the travel operator.
Throughout this subsection $0<q<1$, $D_e=\operatorname{diag}(2q^{1+s})_{s\ge0}$ and
$K=[q^{\max(s,s')}]_{s,s'\ge0}$, so that $\mathcal T=D_eK$ is the travel operator of
{\cite[Definition ``section identity'', \textup{def:model}]{Bonfioli2026bB}}. Put
\[
   \mathbf S:=D_e^{1/2}KD_e^{1/2},\quad \mathbf S[s,s']=2q^{1+(s+s')/2+\max(s,s')},\quad
   F(\lambda,q):=\sum_{k\ge0}\frac{(-2\lambda(1-q))^k\,q^{k^2}}{(q;q)_{2k}}\ \ (\lambda\in\C).
\]
The spectral parameter $\lambda$ is not to be confused with the junction vectors
$\lambda^{(y)}$ of \S\ref{sec:modelassembly}.

\begin{theorem}[the travel determinant]\label{thm:fredholm}
Let $0<q<1$ and $Z=\sqrt{2(1-q)/q}$.
\begin{enumerate}[label=\textup{(\roman*)},leftmargin=2.2em,itemsep=0.1em,topsep=0.2em]
\item $\mathbf S$ is a positive definite trace-class operator on $\ell^2(\Z_{\ge0})$, with
$\operatorname{tr}\mathbf S=2q/(1-q^2)$.
\item $\mathcal T$ and $\mathbf S$ have the same principal minors,
$\det\mathcal T[I,I]=\det\mathbf S[I,I]$ for every finite $I\subset\Z_{\ge0}$, and the same non-zero
eigenvalues, the eigenvectors corresponding by $R=D_e^{1/2}\phi$: if $\mathbf S\phi=\kappa\phi$ with
$\phi\in\ell^2$ and $\kappa\ne0$, then $R\in\ell^1$ and $\mathcal TR=\kappa R$; conversely, if
$R\in\ell^1$ and $\mathcal TR=\kappa R$ with $\kappa\ne0$, then $\phi=D_e^{-1/2}R\in\ell^2$ and
$\mathbf S\phi=\kappa\phi$.
\item For every $\lambda\in\C$, with absolutely convergent series,
\[
   \det(I-\lambda\mathbf S)=\sum_{k\ge0}(-\lambda)^k\!\!\sum_{|I|=k}\det\mathbf S[I,I]
   =F(\lambda,q),
\]
\[
   F(\lambda,q)=\cos(\sqrt\lambda\,Z;q^2)
   =c_q\,(\sqrt\lambda\,Z)^{1/2}J^{(3)}_{-1/2}(\sqrt\lambda\,Z;q^2),
\]
the last two expressions being even in $\sqrt\lambda$. We write $\det(I-\lambda\mathcal T)$ for the
same Fredholm series formed with the minors of $\mathcal T$; by \textup{(ii)} it is the same
function.
\item The eigenvalues of $\mathbf S$ are $1/\lambda_k(q)$, $k\ge1$, with
$0<\lambda_1(q)<\lambda_2(q)<\cdots\to\infty$, all simple; $\sum_k1/\lambda_k(q)=2q/(1-q^2)$ and
$F(\lambda,q)=\prod_{k\ge1}\bigl(1-\lambda/\lambda_k(q)\bigr)$.
\item $1-\Sigma_1=F(1,q)=\det(I-\mathcal T)$, and $S_e=F(q,q)=\det(I-q\mathcal T)=\det(I-M_0)$: the
bulk operator $M_0$ of {\cite[Definition ``section identity'', \textup{def:model}]{Bonfioli2026bB}} is $q\mathcal T$ after the index shift $b=s+1$.
\end{enumerate}
\end{theorem}

\begin{proof}
(i) For $x\in\ell^2$ put $X=D_e^{1/2}x$, which lies in $\ell^1$ by Cauchy--Schwarz. Since
$q^{\max(s,s')}=(1-q)\sum_{t\ge\max(s,s')}q^t$,
\[
   \langle x,\mathbf Sx\rangle=\langle X,KX\rangle
   =(1-q)\sum_{t\ge0}q^t\Bigl|\sum_{s\le t}X_s\Bigr|^2\ \ge\ 0,
\]
the double sums converging absolutely, with equality only if every partial sum of $X$ vanishes,
that is $X=0$, that is $x=0$. $\mathbf S$ is Hilbert--Schmidt, hence bounded, and it is positive
definite; a positive operator with summable diagonal is trace class, with trace
$\sum_s2q^{1+s}q^s=2q/(1-q^2)$.

(ii) $\mathcal T[I,I]=D_e[I,I]\,K[I,I]$ and $\mathbf S[I,I]=D_e[I,I]^{1/2}K[I,I]D_e[I,I]^{1/2}$
with $D_e[I,I]$ diagonal, so the two determinants agree. If $\mathbf S\phi=\kappa\phi$, then
$R=D_e^{1/2}\phi\in\ell^1$ and $\mathcal TR=D_e^{1/2}\mathbf S\phi=\kappa R$. Conversely,
$\mathcal TR=\kappa R$ gives $R=\kappa^{-1}D_e(KR)$ with $|(KR)_s|\le\|R\|_{\ell^1}$, so
$\phi=D_e^{-1/2}R=\kappa^{-1}D_e^{1/2}KR\in\ell^2$ and
$\mathbf S\phi=D_e^{1/2}KR=\kappa\phi$. The operators are not boundedly similar, $D_e^{-1/2}$ being
unbounded; only \textup{(ii)} is used.

(iii) For trace-class $A$, $\det(I-\lambda A)=\sum_k(-\lambda)^k\operatorname{tr}\Lambda^k(A)$
\cite[Ch.~3]{Simon} (the von Koch form; compare Fredholm's formula for continuous kernels,
\cite[Thm.~3.10]{Simon}), and computing the trace of $\Lambda^k(\mathbf
S)$ in the orthonormal basis $e_{s_1}\wedge\cdots\wedge e_{s_k}$, $s_1<\cdots<s_k$, whose diagonal
entries are the principal minors, gives the first equality. The principal minors of the positive
$\mathbf S$ are $\ge0$, so the inner sums converge absolutely, to the values computed next, and the
outer series then converges absolutely for every $\lambda$. For $I=\{s_1<\cdots<s_k\}$ the matrix $K[I,I]$ is $(a_{\max(i,j)})$ with
$a_i=q^{s_i}$ decreasing; subtracting row $i+1$ from row $i$ for $i<k$ leaves a lower triangular
matrix with diagonal $a_1-a_2,\dots,a_{k-1}-a_k,a_k$, so
$\det K[I,I]=q^{s_k}\prod_{i<k}(q^{s_i}-q^{s_{i+1}})$ and, with the gaps
$g_i=s_{i+1}-s_i\ge1$,
\[
   \det\mathbf S[I,I]=2^kq^k\,q^{2(s_1+\cdots+s_k)}\prod_{i<k}\bigl(1-q^{g_i}\bigr).
\]
Since $s_1+\cdots+s_k=ks_1+\sum_{j<k}(k-j)g_j$, the sum over $s_1\ge0$ and $g_1,\dots,g_{k-1}\ge1$
factorises; with $\sum_{g\ge1}q^{2rg}(1-q^g)=q^{2r}(1-q)/\bigl((1-q^{2r})(1-q^{2r+1})\bigr)$,
\[
   \sum_{|I|=k}\det\mathbf S[I,I]
   =\frac{(2q)^k}{1-q^{2k}}\prod_{r=1}^{k-1}\frac{q^{2r}(1-q)}{(1-q^{2r})(1-q^{2r+1})}
   =\frac{(2(1-q))^k\,q^{k^2}}{(q;q)_{2k}} .
\]
Hence $\det(I-\lambda\mathbf S)=F(\lambda,q)$; the case $k=1$ is the trace. Finally
$(2\lambda(1-q))^kq^{k^2}=q^{k^2+k}(\sqrt\lambda\,Z)^{2k}$, since $Z^2=2(1-q)/q$, so
$F(\lambda,q)=\cos(\sqrt\lambda\,Z;q^2)$, and the Hahn--Exton form is the normalisation after
Proposition~\ref{prop:qtrig}.

(iv) For positive trace-class $\mathbf S$ the determinant is the product
$\prod_k(1-\lambda\kappa_k)$ over its eigenvalues $\kappa_k>0$ counted with multiplicity
\cite[Ch.~3]{Simon}, so the $1/\kappa_k$ are the zeros of $F(\cdot,q)$, the multiplicity of
$\kappa_k$ being the order of the zero. By \textup{(iii)} these are the $\lambda=x^2/Z^2$ with $x>0$
a zero of $\cos(x;q^2)=c_qx^{1/2}J^{(3)}_{-1/2}(x;q^2)$ ($\cos(\cdot;q^2)$ is even with
$\cos(0;q^2)=1$). For $\nu>-1$ the zeros of $J^{(3)}_\nu(\cdot;q^2)$ are real and simple and form
an infinite sequence \cite[Cor.~3.2, Lemma~3.3, Thm.~3.4]{KoelinkSwarttouw1994}; here
$\nu=-\tfrac12$ with base $q^2$. So the eigenvalues are simple and infinitely many, and the sum of
the $1/\lambda_k$ is the trace.

(v) At $\lambda=1$, $F(1,q)=\sum_k(-1)^k(2(1-q))^kq^{k^2}/(q;q)_{2k}=1-\Sigma_1$
(Proposition~\ref{prop:qtrig}); at $\lambda=q$, $F(q,q)=\sum_k(-2(1-q))^kq^{k^2+k}/(q;q)_{2k}=S_e$
by \eqref{eq:blocks}. And $M_0[s+1,s'+1]=2q^{s+1}q^{\max(s,s')+1}=q\,\mathcal T[s,s']$, so $M_0$
has the principal minors of $q\mathcal T$.
\end{proof}

\noindent The identity $\det(I-\lambda\mathbf S)=F(\lambda,q)$ was checked in double precision
against truncated matrices at $q=0.3,0.6,0.8$ and four values of $\lambda$ (one of them non-real),
to within $3.2\cdot10^{-15}$, and the gap sum against a brute-force sum of minors for $k\le3$ at
$q=0.3$.

\begin{remark}[a $q$-string]\label{rem:qstring}
With the Jackson integral $\int_0^1f\,d_qt=(1-q)\sum_{s\ge0}q^sf(q^s)$, let
$K_qf(x)=\int_0^1\min(x,t)f(t)\,d_qt$ on the lattice $x\in\{q^s\}$. At $x=q^s$,
$(KD_eu)_s=\sum_{s'}q^{\max(s,s')}2q^{1+s'}u_{s'}=\tfrac{2q}{1-q}(K_qu)(q^s)$, that is,
\[
   \mathcal T^{\!\top}=KD_e=\frac{2q}{1-q}\,K_q ,
\]
so up to the transposition, which changes neither the principal minors nor the eigenvalues,
$\mathcal T=\tfrac{2q}{1-q}K_q$. The kernel $\min(x,t)$ is the Green function of $-d^2/dx^2$ on
$[0,1]$ with a Dirichlet end at $0$ and a Neumann end at $1$: $\mathcal T$ is the compliance operator
of a $q$-string with uniform Jackson mass, fixed at $0$ and free at $1$, and $u=KR$ is its
displacement ($u_s\to0$ as $x=q^s\to0$). As $q\to1$ the kernel tends to the classical one, with
eigenvalues $4/((2k-1)^2\pi^2)$; this is the heuristic behind
$\lambda_k(q)\sim(1-q)(2k-1)^2\pi^2/8$, proved for fixed $k$ in
Proposition~\ref{prop:infpolesalt}.
\end{remark}

\begin{corollary}[the travel poles are Hahn--Exton zeros]\label{cor:qbessel}
A point $q\in(0,1)$ is a travel pole if and only if
\[
   J^{(3)}_{-1/2}\bigl(Z(q);q^2\bigr)=0,\qquad Z(q)=\sqrt{2(1-q)/q},
\]
if and only if $1$ is an eigenvalue of $\mathcal T(q)$, that is $\lambda_k(q)=1$ for some $k$. That
$k$ is unique, the eigenvalue is simple, its eigenvector is the travel null vector $R$ of
Proposition~\ref{prop:shape}, and $\phi=D_e^{-1/2}R$ is the corresponding eigenvector of $\mathbf S$.
\end{corollary}

\begin{proof}
Theorem~\ref{thm:fredholm}\textup{(iii)--(v)}, with $c_q>0$ and $Z>0$, and
Proposition~\ref{prop:qtrig}; uniqueness of $k$ is the simplicity in
Theorem~\ref{thm:fredholm}\textup{(iv)}, and the eigenvector is Theorem~\ref{thm:fredholm}\textup{(ii)}
with {\cite[Lemma ``the null space is at most one-dimensional'', \textup{lem:nullspace}]{Bonfioli2026bB}}.
\end{proof}

\noindent The zero theory used is that of Koelink--Swarttouw \cite{KoelinkSwarttouw1994}, at
$\nu=-\tfrac12$ and base $q^2$, where its hypothesis $\nu>-1$ holds. The eigenvalues also agree with
the Jacobi-matrix description of the zeros of $J^{(3)}_\nu$ by \v{S}tampach and \v{S}\v{t}ov\'{\i}\v{c}ek
\cite{StampachStovicek2014}, in its Friedrichs extension; the positive form of $\mathbf S$ is the
Friedrichs realisation, on which $u_s=(KR)_s\to0$. At $q_1,\dots,q_{10}$, located at $90$-digit
precision, $|J^{(3)}_{-1/2}(Z;q^2)|\le2\cdot10^{-79}$, and exactly $m-1$ zeros of
$\cos(\cdot;q_m^2)$ lie below $Z(q_m)$ (a numerical check, not a certificate).

\begin{proposition}[infinitely many travel poles, without $T_2$]\label{prop:infpolesalt}
\begin{enumerate}[label=\textup{(\alph*)},leftmargin=2.2em,itemsep=0.1em,topsep=0.2em]
\item $\lambda_1(q)>(1-q^2)/(2q)$ on $(0,1)$. In particular every $\lambda_k(q)>1$ for
$q\le\sqrt2-1$, and there is no travel pole in $(0,\sqrt2-1]$.
\item For each fixed $k$, $\lambda_k(q)/(1-q)\to(2k-1)^2\pi^2/8$ as $q\to1^-$.
\item Each $\lambda_k$ is continuous on $(0,1)$, and there is $c_k\in(\sqrt2-1,1)$ with
$\lambda_k(c_k)=1$. The $c_k$ are pairwise distinct travel poles; in particular there are infinitely
many travel poles.
\end{enumerate}
None of this uses $T_2$, {\cite[Lemma ``absolute bound on $T_2$'', \textup{lem:T2abs}]{Bonfioli2026bB}}, or a single-crossing property.
\end{proposition}

\begin{proof}
(a) The largest eigenvalue of the positive trace-class $\mathbf S$ is strictly less than its trace,
since $\mathbf S$ has infinitely many positive eigenvalues: $1/\lambda_1<2q/(1-q^2)$. And
$(1-q^2)/(2q)\ge1$ exactly for $q\le\sqrt2-1$. At such $q$ no $\lambda_k$ equals $1$, so
$1-\Sigma_1=F(1,q)\ne0$ by Theorem~\ref{thm:fredholm}\textup{(iv)}.

(b) Put $\varepsilon=1-q$ and $\lambda=\varepsilon y$. The $k$-th term of $F(\varepsilon y,q)$ is
$t_k(y,\varepsilon)=(-1)^k(2y\varepsilon^2)^kq^{k^2}/(q;q)_{2k}$. By the arithmetic--geometric mean
inequality, $1-q^n=\varepsilon\sum_{i<n}q^i\ge n\varepsilon q^{(n-1)/2}$, so
$(q;q)_{2k}\ge\varepsilon^{2k}(2k)!\,q^{k(2k-1)/2}$ and
\[
   |t_k(y,\varepsilon)|\ \le\ \frac{(2|y|)^k\,q^{k/2}}{(2k)!}\ \le\ \frac{(2|y|)^k}{(2k-1)!!},
\]
uniformly in $\varepsilon\in(0,1)$. Each $t_k(y,\varepsilon)$ tends to $(-1)^k(2y)^k/(2k)!$ as
$\varepsilon\to0$, because $q^{k^2}\to1$ and $(q;q)_{2k}/\varepsilon^{2k}\to(2k)!$. The domination
is summable locally uniformly in $y\in\C$, so $F(\varepsilon y,1-\varepsilon)\to\cos\sqrt{2y}$
uniformly on compact sets. The limit has the simple zeros $y_j=(2j-1)^2\pi^2/8$, and the zeros of
$y\mapsto F(\varepsilon y,1-\varepsilon)$ are the $\lambda_j(q)/\varepsilon$, positive and ordered.
By Hurwitz's theorem, for $\varepsilon$ small each disc $|y-y_j|<\delta$, $j\le k$, contains exactly
one of them and $|y|\le y_k+\delta$ contains no other; hence $\lambda_j(q)/\varepsilon\to y_j$ for
$j\le k$.

(c) $q\mapsto\mathbf S(q)$ is continuous in trace norm, hence in operator norm, on $(0,1)$: its
entries are continuous and $\sum_{s,s'}|\mathbf S[s,s']|$ converges locally uniformly. The ordered
eigenvalues of positive compact operators are $1$-Lipschitz in operator norm (min--max), so each
$1/\lambda_k$, and so each $\lambda_k$, is continuous. By (a) $\lambda_k>1$ on $(0,\sqrt2-1]$, and by
(b) $\lambda_k(q)\to0$ as $q\to1$, so the intermediate value theorem gives $c_k$; it is a travel pole
by Corollary~\ref{cor:qbessel}. If $c_j=c_k$ with $j\ne k$, then $\lambda_j(c_k)=\lambda_k(c_k)=1$,
against simplicity (Theorem~\ref{thm:fredholm}\textup{(iv)}).

The same argument with $\lambda_k(q)/q$ in place of $\lambda_k(q)$ gives {\cite[Lemma ``infinitude of the zeros of the bulk block'', \textup{lem:infzeros}]{Bonfioli2026bB}}
without $T_2$ too: $S_e=F(q,q)=\prod_k(1-q/\lambda_k(q))$ vanishes exactly where some
$\lambda_k(q)=q$; $\lambda_k(q)/q>(1-q^2)/(2q^2)>1$ for $q<1/\sqrt3$ by (a); $\lambda_k(q)/q\to0$ by
(b); and two branches cannot both equal $q$ at the same point.
\end{proof}

\subsection{Local finiteness, and accumulation at $q=1$}\label{sec:polefinite}

\begin{lemma}[The travel poles are locally finite]\label{lem:polefinite}
For every $\tau_\star>0$ there are only finitely many travel poles with $\tau_m>\tau_\star$.
\end{lemma}

\begin{proof}
By Proposition~\ref{prop:qtrig} the travel poles are the zeros in $q\in(0,1)$ of
$g(q):=\cos(Z;q^2)=C(Z^2)$ with $Z^2=2(1-q)/q$, i.e.\ of
\[
 g(q)=\sum_{j\ge0}(-1)^j\,\frac{2^j\,q^{j^2}(1-q)^j}{(q;q)_{2j}} .
\]
First, $g$ has no zero near $q=0$. For $0<q\le\tfrac1{10}$ one has
$(q;q)_{2j}\ge\prod_{i\ge1}(1-q^i)\ge1-\tfrac q{1-q}\ge\tfrac89$ and $(1-q)^j\le1$, so
$|g(q)-1|\le\tfrac98\sum_{j\ge1}2^jq^{j^2}\le\tfrac98\bigl(2q+4q^4/(1-q)\bigr)\le0.24<1$, whence
$g(q)>0$ there.

Second, $g$ is real-analytic on $(0,1)$: the series converges locally uniformly, since on
$q\le1-\delta$ the $j$-th term is bounded by $2^jq^{j^2}/(q;q)_\infty$ and $q^{j^2}$ dominates $2^j$.
Hence the zeros of $g$ are isolated in $(0,1)$, and any zero set with an accumulation point in the
\emph{open} interval would force $g\equiv0$, contradicting $g(\tfrac1{10})>0$.

Now fix $\tau_\star>0$. The poles with $\tau_m>\tau_\star$ lie in $q\in[\tfrac1{10},e^{-\tau_\star}]$,
a compact subinterval of $(0,1)$. An analytic function on a domain whose zeros are isolated has only
finitely many zeros on any compact subset; hence there are finitely many such poles.
\end{proof}

\begin{remark}
Only local finiteness is proved here; that the poles \emph{do} accumulate at $q=1$
(equivalently $\tau_m\downarrow0$) needs in addition that there are infinitely many of them, which
is Proposition~\ref{prop:infpolesalt}\textup{(c)} (see the paragraph after this remark), or,
localised in windows of $w$, \cite[Lemma ``infinitude and accumulation of the travel poles'',
\textup{lem:infpoles}]{Bonfioli2026bB}. The present lemma says merely that
they cannot accumulate anywhere inside $(0,1)$, so that any threshold $\tau_\star$ leaves a finite
set to be checked directly. Its proof also shows $g>0$ on $(0,\tfrac1{10}]$, so every travel pole
has $q_m>\tfrac1{10}$. Besides the accumulation statement used in Theorems~\ref{thm:travelres},
\ref{thm:V} and~\ref{thm:U}, it makes the splice in
\cite[Proposition \textup{prop:selower}]{Bonfioli2026bB} well posed and the exceptional set of
\cite[Theorem ``the amplitude estimate $(\star)$'', \textup{thm:star}]{Bonfioli2026bB} finite, and it
gives the accumulation clause of \cite[Lemma \textup{lem:infpoles}]{Bonfioli2026bB}.
\end{remark}

\noindent With Proposition~\ref{prop:infpolesalt}\textup{(c)}, Lemma~\ref{lem:polefinite} gives:
\emph{the travel poles form an infinite subset of $(\tfrac1{10},1)$ whose only accumulation point
is $q=1$}; in particular, for every $\tau_\star>0$ all but finitely many travel poles have
$\tau_m\le\tau_\star$. This is the form in which the pole set is used below. It needs no estimate
on $T_2$ and no numerical input.

\subsection{The numerators, and why they do not vanish at the poles}\label{sec:numerators}

Proposition~\ref{prop:qtrig} identifies the two \emph{denominators}. The two numerators are the
companion $q$-sine at the same two arguments, and that is what settles their non-vanishing at the
travel poles. Throughout, $\mathfrak s(u)=q^{-1/2}\sin(q^{1/2}u;q^2)$ is the half-step $q$-sine of
{\cite[Lemma ``the gate block in closed form'', \textup{lem:P12closed}]{Bonfioli2026bB}}, $\sin(z;q^2)=\sum_{j\ge0}(-1)^jq^{j^2+j}z^{2j+1}/(q;q)_{2j+1}$.

\begin{lemma}[the numerators in closed form]\label{lem:qsinenum}
Identically in $0<q<1$,
\[
   \Sigma_0=\frac{2q}{Z}\,\mathfrak s(Z),\qquad S_0=\frac{2q}{z_0}\,\mathfrak s(z_0).
\]
With Proposition~\ref{prop:qtrig} this identifies all four blocks as one $q$-trigonometric pair
sampled at two arguments a half $q$-step apart: $(\Sigma_0,\,1-\Sigma_1)$ is
$\bigl(\tfrac{2q}u\mathfrak s(u),\,\cos(u;q^2)\bigr)$ at $u=Z$, and $(S_0,\,S_e)$ is the same pair
at $u=z_0=\sqrt q\,Z$. In particular $S_0=\tfrac{2q}{1-q}S_o$ with $S_o$ the odd bulk block of
{\cite[Lemma ``the gate block in closed form'', \textup{lem:P12closed}]{Bonfioli2026bB}}, which is the second identity of the dictionary {\cite[\textup{(eq:dict)}]{Bonfioli2026bB}} read as a
statement about the explicit series, and is therefore independent of the transfer model.
\end{lemma}

\begin{proof}
Collapse the $k$-recursion \eqref{eq:krec} on the even ladders, exactly as
Proposition~\ref{prop:qtrig} does on the odd ones. From $\alpha_k=2q^{k+1}/(1-q^{k+1})$,
\[
  |\gamma_k|=\alpha_k-\alpha_{k+1}=\frac{2(1-q)q^{k+1}}{(1-q^{k+1})(1-q^{k+2})},\qquad
  |C_k|=\frac{2(1-q)q^{k+2}}{(1-q^{k+1})(1-q^{k+2})} ,
\]
$\gamma_k$ and $C_k$ themselves being negative on $(0,1)$, so both even ladders alternate. For the bulk one,
$\sum_{i<j}(2i+1)=j^2$ and the denominators telescope to $(q;q)_{2j}$, giving
$\prod_{i<j}|\gamma_{2i}|=[2(1-q)]^jq^{j^2}/(q;q)_{2j}$ and hence
\[
  S_0=\sum_{j\ge0}(-1)^j\alpha_{2j}\prod_{i<j}|\gamma_{2i}|
     =2\sum_{j\ge0}(-1)^j\frac{[2(1-q)]^j\,q^{(j+1)^2}}{(q;q)_{2j+1}} .
\]
Write $q^{(j+1)^2}=q\,q^{j^2+j}q^{j}$ and $[2(1-q)]^jq^{j}=\bigl(q^{1/2}z_0\bigr)^{2j}$: the sum is
$(2q/(q^{1/2}z_0))\sin(q^{1/2}z_0;q^2)=(2q/z_0)\mathfrak s(z_0)$. For the travel ladder,
$\sum_{i<j}(2i+2)=j^2+j$ and the denominators telescope to $(q;q)_{2j}$ again, so
$\prod_{i<j}|C_{2i}|=[2(1-q)]^jq^{j^2+j}/(q;q)_{2j}$ and, with $A_{2j}=2q/(1-q^{2j+1})$,
\[
  \Sigma_0=2q\sum_{j\ge0}(-1)^j\frac{[2(1-q)]^j\,q^{j^2+j}}{(q;q)_{2j+1}}
          =\frac{2q}{z_0}\sin(z_0;q^2)=\frac{2q^{3/2}}{z_0}\,\mathfrak s(Z)=\frac{2q}{Z}\,\mathfrak s(Z),
\]
using $[2(1-q)]^j=z_0^{2j}$ and $z_0=q^{1/2}Z$. The $S_o$ statement is the identity $S_o=\tfrac{z_0}2\mathfrak s(z_0)$ of
\cite[Lemma ``the gate block in closed form'', \textup{lem:P12closed}]{Bonfioli2026bB} combined with $z_0^2=2(1-q)$. Both collapses are the exponent
bookkeeping machine-checked in \texttt{QTrigIdentities.lean} and are verified numerically off the
poles, hence as identities, in \texttt{qnumerator\_certificate.py}.
\end{proof}

\begin{lemma}[half-step rules and the pole invariant]\label{lem:halfstep}
Identically in $0<q<1$, with $c(u)=\cos(u;q^2)$:
\begin{enumerate}[label=\textup{(\roman*)},leftmargin=2.2em,itemsep=0.1em,topsep=0.2em]
\item $c(qu)=c(u)+q^{3/2}u\,\mathfrak s(q^{1/2}u)$;
\item $\mathfrak s(qu)=\mathfrak s(u)-u\,c(q^{1/2}u)$;
\item $c(u)\,c(q^{1/2}u)+q^{3/2}\,\mathfrak s(u)\,\mathfrak s(q^{1/2}u)=1$.
\end{enumerate}
Consequently, at every travel pole, where $c(Z)=0$:
\[
   q^{3/2}\,\mathfrak s(z_0)\,\mathfrak s(Z)=1 ,
\]
so neither $\mathfrak s(Z)$ nor $\mathfrak s(z_0)$ vanishes there.
\end{lemma}

\begin{proof}
(i) and (ii) are termwise. From $(q;q)_{2j}=(q;q)_{2j-1}(1-q^{2j})$,
\[
   c(u)-c(qu)=\sum_{j\ge1}(-1)^jq^{j^2+j}u^{2j}/(q;q)_{2j-1};
\]
shifting $j\mapsto j+1$ and factoring
out $-qu$ leaves $-qu\sin(qu;q^2)=-q^{3/2}u\,\mathfrak s(q^{1/2}u)$, which is (i). From
$(q;q)_{2j+1}=(q;q)_{2j}(1-q^{2j+1})$,
$\mathfrak s(u)-\mathfrak s(qu)=q^{-1/2}\sum_{j\ge0}(-1)^jq^{j^2+j}(q^{1/2}u)^{2j+1}/(q;q)_{2j}$;
factoring out $q^{1/2}u$ leaves $u\,c(q^{1/2}u)$, which is (ii). This is the exponent
matching machine-checked as \texttt{cshift\_exponent} and \texttt{sshift\_exponent}. For (iii),
substitute (i) and (ii) into the Hahn--Exton $q$-Wronskian
$\mathfrak s(u)c(qu)-c(u)\mathfrak s(qu)=u$ (\cite[eq.~(4.3.8)]{Swarttouw1992};
\cite[eq.~(4.21)]{Daalhuis1994}): the $\mathfrak s(u)c(u)$ terms cancel and what is left is
$u\bigl[c(u)c(q^{1/2}u)+q^{3/2}\mathfrak s(u)\mathfrak s(q^{1/2}u)\bigr]$, so the $q$-Wronskian and
(iii) are equivalent for $u\ne0$. Identity (iii) is the $q$-Pythagorean identity of
\cite{KS1992}. Taking $u=Z$ and $c(Z)=0$ gives the pole invariant, which is the step already used
inside the proof of {\cite[Proposition \textup{prop:P12second}]{Bonfioli2026bB}}.
\end{proof}

\begin{proposition}[the numerators do not vanish at the travel poles]\label{prop:numnonvanish}
At every travel pole $q_m$,
\begin{equation}\label{eq:numprod}
   \Sigma_0(q_m)\,S_0(q_m)\;=\;\frac{2q_m}{1-q_m}\;\ne\;0 .
\end{equation}
In particular $\Sigma_0(q_m)\ne0$ and $S_0(q_m)\ne0$, at \emph{every} travel pole, with no
asymptotic input and no exceptional set. Moreover $|\Sigma_0(q_m)|\,|S_0(q_m)|\ge q_mw_m^2$, so at
least one of the two numerators exceeds $\sqrt{q_m}\,w_m$ in modulus.
\end{proposition}

\begin{proof}
By Lemma~\ref{lem:qsinenum}, $\Sigma_0S_0=\bigl(4q^2/(z_0Z)\bigr)\mathfrak s(Z)\mathfrak s(z_0)$,
and by Lemma~\ref{lem:halfstep} the product $\mathfrak s(Z)\mathfrak s(z_0)$ equals $q^{-3/2}$ at a
travel pole. Since $z_0Z=z_0^2/\sqrt q=2(1-q)/\sqrt q$, the two powers of $q$ combine to
$4q^2\cdot q^{-3/2}\cdot\sqrt q/(2(1-q))=2q/(1-q)$. The size statement follows from
$1-q\le\tau=2/w^2$.
\end{proof}

\section{The transfer model and the local form at a travel pole}\label{sec:modelassembly}

This section connects the explicit operators to the growth series. We first collect the objects
the transfer model is phrased in; their construction, and every identity between them that is used
below, is proved in \cite[\S\,\textup{sec:assembly}]{Bonfioli2026bB}.

\subsection{Notation}\label{sec:notation}

Throughout, $|q|<1$, $x=q^{1/2}$, and at a travel pole $q_m$ we write $\tau_m=-\log q_m$,
$w_m=\sqrt{2/\tau_m}$ and $x_m=\sqrt{q_m}$.

\emph{The two operators.} On $\ell^1(\Z_{\ge0})$ the \emph{travel operator} is $\mathcal T=D_eK$
with $K=[q^{\max(s,s')}]_{s,s'\ge0}$ and $D_e=\operatorname{diag}(e_s)$, $e_s=2q^{1+s}$; on
$\ell^1(\Z_{\ge1})$ the \emph{bulk operator} is $M_0[b,a]=2q^b\,q^{\max(a,b)}$. These are the travel
and bulk instances of the section identity of
\cite[Definition ``section identity'', \textup{def:model}]{Bonfioli2026bB}; both are compact, and
$(I-\mathcal T)^{-1}$ and $R_0=(I-M_0)^{-1}$ are meromorphic on $|q|<1$ with poles only where
$1-\Sigma_1=0$, respectively $S_e=0$
\cite[Corollary ``singularities of the two gapless resolvents'', \textup{cor:sing}]{Bonfioli2026bB}.
Put $E_b=2q^b$, $u_b=2q^{2b}$, $v_a=q^a$ and
\[
   t_1=v^{\!\top}R_0u ,
\]
an explicit bulk quantity; one has $t_1=P_{12}/S_e$ identically, where $P_{12}$ is the off-diagonal
entry of the bulk cocycle, a Hahn--Exton function $J^{(3)}_{3/2}$ in closed form
\cite[Corollary ``the cocycle is the same object'', \textup{cor:cocycle}]{Bonfioli2026bB}.

\emph{The gap term.} With $\mathfrak g=\mathfrak g(q,y)=\sum_{L\ge1}q^Ly^{L-1}=q/(1-qy)$, the
gap-marked bulk system is
\begin{equation}\label{eq:gapkernel}
   K_{\mathfrak g}(a,b)=q^{\max(a,b)}+q^{a+b}\,\mathfrak g,\qquad
   P_b=2q^b\Bigl(q^b+\sum_{a\ge1}K_{\mathfrak g}(a,b)P_a\Bigr)\quad(b\ge1).
\end{equation}
The gap term is of rank one, $2q^b\cdot q^{a+b}=u_bv_a$, and by the Sherman--Morrison formula
$P=B_0\,R_0u$ with
\[
   B_0=1+\mathfrak g\sum_{b\ge1}q^bP_b=\frac1{1-\mathfrak g\,t_1},\qquad
   B_{0z}=1+y\,\mathfrak g\sum_{b\ge1}q^bP_b
\]
\cite[Proposition ``M\"obius factorisation'', \textup{prop:mobius}; Remark ``the bulk sum of the
model'', \textup{rem:bulksource}]{Bonfioli2026bB}. The two evaluations of the gap kernel are
$\mathfrak g_V=\mathfrak g(q,1)=\tfrac q{1-q}$ and $\mathfrak g_U=\mathfrak g(q,q)=\tfrac q{1-q^2}$.

\emph{The junction vectors.} For $s\ge0$,
\begin{equation}\label{eq:junctionsym}
   \mu^{(y)}_s=B_0\,x^{2s+1}+\sum_{b\ge1}\frac{P_b}{2}
   \Bigl(x^{\max(2b-1,\,2s+1)}+x^{\max(2b+1,\,2s+1)}\Bigr),
\end{equation}
\begin{equation}\label{eq:farjunction}
   E^{(y)}_s=B_0^{[s]}x^{2s}+B_0\,x^{2s+2}+\sum_{b\ge1}P_b
   \Bigl(x^{\max(2s,\,2b)}+x^{\max(2s+2,\,2b)}\Bigr),\qquad
   B_0^{[0]}=B_{0z},\ B_0^{[s]}=B_0\ (s\ge1),
\end{equation}
and $\lambda^{(y)}=E^{(y)}+2\mu^{(y)}$. These are the near- and far-junction vectors of the marker
sites; the site costs they encode are computed in
\cite[Corollary ``the marker junctions'', \textup{cor:marker}]{Bonfioli2026bB}, and the rank of the
junction is discussed in \cite[Proposition ``the marker junction is not rank one'',
\textup{prop:junction}]{Bonfioli2026bB}.

\emph{The bivariate series.} With $c(g)$ the cycle count and $\ell_R(g)$ the relaxed length,
\[
   \mathcal W(x,y)=\sum_{n,c\ge0}w_{n,c}\,x^ny^c,\qquad
   w_{n,c}=\#\{g:\ell_R(g)=n,\ c(g)=c\},
\]
so that $V(x)=\mathcal W(x,1)$ and, since $\ell_T=\ell_R+2c$ for every element (formalised in Lean,
\texttt{metricAll}), $U(x)=\mathcal W(x,x^2)=\mathcal W(x,q)$
\cite[\S\,\textup{sec:bridge}]{Bonfioli2026bB}.

\emph{The gate.} Write $s=\mathfrak g_Vt_1=\tfrac q{1-q}t_1$ and $b_0=S_0/(1-S_1)$. The gate is
\begin{equation}\label{eq:gate}
   b_0>0\qquad\text{and}\qquad s<1 ;
\end{equation}
it holds at every travel pole, with $|s|\le0.983$ for $\tau_m\le5\cdot10^{-3}$ and by direct
evaluation at the six poles above that threshold, whose enumeration has the standing stated in
\cite[Theorem ``the amplitude estimate $(\star)$'', \textup{thm:star}]{Bonfioli2026bB}
\textup{(}see \cite[Corollary ``the gate closes'', \textup{app:gate}]{Bonfioli2026bB}\textup{)}. The half $|s|<1$
holds unconditionally, and the half $b_0>0$ through the closed form of $b_0$, which rests on the
transfer model.

\subsection{The model \textup{(M)}}

The strand-walk input that identifies the operators above with the growth series is stated once and
named.

\begin{theorem}[the transfer model \textup{(M)}]\label{hyp:model}\label{thm:model}
\textup{(M1)} \emph{Local cost law.} The site cost between adjacent active edges carrying deposits
$d,d'$ is $x^{\max(|d|,|d'|)}$, independent of the crossing counts and of the travel indicator.
Bulk deposits are even, $|d|=2s$, and travel deposits odd, $|d|=2s+1$, so this cost is
$q^{\max(s,s')}$ in the bulk and $x^{2\max(s,s')+1}$ on the travel side, and the edge weights are
$2q^s$ and $2x^{2s+1}$ respectively: {\cite[Definition ``section identity'', \textup{def:model}]{Bonfioli2026bB}} in its two instances.
At the near marker site the cost is $x^{\mathrm{Site}_0(d_L,d_R)}$ with
$\mathrm{Site}_0(d_L,d_R)=\max(|d_L-1|,|d_R|)$, $d_L$ the last bulk deposit and $d_R$ the first
travel deposit, for all four marker data $(\varepsilon^\ast,\delta^\ast)$; at the far marker site
it is the $(\varepsilon^\ast,\delta^\ast)$-dependent cost of {\cite[Corollary ``the marker junctions'', \textup{cor:marker}]{Bonfioli2026bB}}.
\textup{(M2)} \emph{Shield law.} A maximal interior run of $L$ forced gap edges costs $q^L$ and
carries $L-1$ isolated cycles; hence the gap-bridge kernel is $\mathfrak g(q,y)=q/(1-qy)$ and the
gap-marked bulk kernel is the $K_{\mathfrak g}$ of \eqref{eq:gapkernel}.
\textup{(M3)} \emph{Decomposition.} The bivariate series of {\cite[\S\,\textup{sec:bridge}]{Bonfioli2026bB}} is
$\mathcal W=\mathcal W_++\mathcal W_-+\mathcal W_0$, the sectors $k^\ast>0$, $k^\ast<0$, $k^\ast=0$,
with
\begin{equation}\label{eq:assembly}
 \begin{gathered}
   \mathcal W_+=x^{-1}\bigl\langle\mu^{(y)},(I-\mathcal T)^{-1}D_e\lambda^{(y)}\bigr\rangle,\qquad
   \mathcal W_-=x^{-1}\bigl\langle\tfrac12E^{(y)},(I-\mathcal T)^{-1}D_e\lambda^{(y)}\bigr\rangle,\\
   \mathcal W_++\mathcal W_-=\frac1{2x}\bigl\langle\lambda^{(y)},(I-\mathcal T)^{-1}D_e\lambda^{(y)}\bigr\rangle,
 \end{gathered}
\end{equation}
where $\mathcal T=D_eK$ is the travel operator of {\cite[Definition ``section identity'', \textup{def:model}]{Bonfioli2026bB}},
$D_e=\operatorname{diag}(e_s)$ with $e_s=2q^{1+s}$, and $\mu^{(y)},E^{(y)},\lambda^{(y)}$ are as in
\eqref{eq:junctionsym}--\eqref{eq:farjunction}. The sum over the marker data
$(\varepsilon^\ast,\delta^\ast)$ is already carried out in these vectors: $\delta^\ast=1$ at the far
site gives $2\mu$ and $\delta^\ast=0$ gives $E$ when $k^\ast>0$, and the roles are exchanged when
$k^\ast<0$. The $k^\ast=0$ sector is
$\mathcal W_0=A_0+B_0A_1+B_0^2A_2$, with $A_i$ built from bulk quantities alone; it contains no
travel resolvent.
\end{theorem}

\begin{proof}[Proof and standing]
Let $(k^\ast,\varepsilon^\ast,\delta^\ast,d)$ parametrise the elements of $W_{\mathrm{univ}}$ as
in \cite[Theorem ``Normal form and lamp model'', \textup{thm:normal-form}]{Bonfioli2026a}. Two facts
are formalised in Lean~4 with the standard axioms only:
\begin{itemize}
\item the affine realisation intertwines the three generators and is injective, and every element
of the model is reachable, so the word length of the model is the Cayley length in
$W_{\mathrm{univ}}$ (\texttt{RJPhi.lean}: \texttt{intertwine\_s1}, \texttt{intertwine\_s2},
\texttt{intertwine\_s3}, \texttt{injective}; reachability is \texttt{reaches\_of\_phiZ} in
\texttt{PhiLipschitz.lean});
\item $\mathrm{wordLength}=\ell_R+2c$ for every element, with $\ell_R$ and $c$ the site-local
closed forms \texttt{lRTrue}, \texttt{cTrue} of \texttt{CorrectedSpan.lean}
(\texttt{RJMetricAll.lean}, theorem \texttt{metricAll}).
\end{itemize}
Given these, $\ell_R$ and $c$ are sums of site and edge costs whose only non-bulk sites are $0$ and
$k^\ast$ (\cite[Corollaries \textup{cor:localcost} and \textup{cor:marker}]{Bonfioli2026bB}; a gap run of length $L$ contributes
$q^Ly^{L-1}$ and a travel-interior site is never a cut). Summing these site costs chain by chain
gives \eqref{eq:gapkernel} with source $u$, the vectors
\eqref{eq:junctionsym}--\eqref{eq:farjunction}, and the three sectors of \eqref{eq:assembly}; this
is {\cite[Theorem ``\textup{(M3$'$)}'', \textup{thm:M3prime}]{Bonfioli2026bB}}, proved in {\cite[Appendix \textup{app:M3prime}]{Bonfioli2026bB}}, and the regrouping of
$\mathcal W_0$ is {\cite[Remark ``the regrouping of $\mathcal W_0$'', \textup{M3:rem:W0regroup}]{Bonfioli2026bB}}. That last step is a hand proof, with two
independent adversarial reviews, and it is not formalised. Every identity holds in
$\Z[y][[x]]$, and the series agree with enumeration as recorded in {\cite[\S\,\textup{sec:bridge}]{Bonfioli2026bB}}.
\end{proof}

\noindent One caveat: $\ell_R$ in $\mathcal W$ is the closed-form relaxed length of
\cite{Bonfioli2026a}. That it also equals the minimum over relaxed realisations
(\cite[Remark \textup{rem:metric-formula-status}]{Bonfioli2026a}) affects only the \emph{interpretation}
of $V$ as a relaxed count, not $U$ and not the series identities. How \textup{(M)} was carried in
earlier versions, and which of its weak links the proof removes, is recorded in
{\cite[\S\,\textup{sec:inputL}]{Bonfioli2026bB}}.

$V=\mathcal W(x,1)$ and $U=\mathcal W(x,q)$ ({\cite[\S\,\textup{sec:bridge}]{Bonfioli2026bB}}), the two evaluations differing
only through $\mathfrak g_V=\tfrac q{1-q}$, $\mathfrak g_U=\tfrac q{1-q^2}$ and the slot $B_{0z}$;
and $\lambda^{(y)},\mu^{(y)},\mathcal W_0$ are holomorphic at a travel pole $q_m$, by
{\cite[\textup{(eq:smresolvent)}]{Bonfioli2026bB}}, as soon as $S_e(q_m)\ne0$ and $1-\mathfrak gt_1\ne0$ there, which is
{\cite[Proposition \textup{prop:selower}]{Bonfioli2026bB}} and {\cite[Corollary ``the gate closes'', \textup{app:gate}]{Bonfioli2026bB}} ($|\mathfrak g_Vt_1|\le0.983$, and
$|\mathfrak g_Ut_1|<|\mathfrak g_Vt_1|$). In particular $\mathcal W_0$ has no travel pole: by
{\cite[Remark ``the regrouping of $\mathcal W_0$'', \textup{M3:rem:W0regroup}]{Bonfioli2026bB}} it is a polynomial in $B_0$ with coefficients built from the bulk
vector $P$ and monomials, so its only possible singularities are at $S_e=0$ and at
$1-\mathfrak gt_1=0$.

The growth series are defined as power series, with radius of convergence
$x_1=\sqrt{q_1}=1/\beta_2=0.6704\ldots$ (Corollary~\ref{cor:beta}), and the local forms below are
statements about functions on a larger disc. The following lemma supplies that continuation and
identifies it with the series.

\begin{lemma}[meromorphic continuation]\label{lem:mero}
Let $y$ be either a complex constant with $|y|\le1$ or $y=x^2$. Then the specialisation
$x\mapsto\mathcal W(x,y)$ of the series of {\cite[Theorem ``\textup{(M3$'$)}'', \textup{thm:M3prime}]{Bonfioli2026bB}} converges near $x=0$ and is
the Taylor series at $0$ of a single-valued meromorphic function on the disc $|x|<1$, given there
by \eqref{eq:assembly} and $\mathcal W_0=A_0+B_0A_1+B_0^2A_2$ read as analytic expressions. Its
poles in $0<|x|<1$ lie over points $q=x^2$ at which $1-\Sigma_1(q)=0$, $S_e(q)=0$ or
$1-\mathfrak g(q,y)\,t_1(q)=0$. In particular $V=\mathcal W(x,1)$ and $U=\mathcal W(x,x^2)$ are
meromorphic on $|x|<1$, and $G^T_0=\Sigma_0/(1-\Sigma_1)$ is meromorphic on $|q|<1$.
\textup{(}The continuation concerns the explicit operators of \cite[\S\S\,\textup{sec:sect}--\textup{sec:mob}]{Bonfioli2026bB}; its
identification with $U$ and $V$ is Theorem~\ref{thm:model}.\textup{)}
\end{lemma}

\begin{proof}
Put $q=x^2$, so $|q|<1$; $\mathfrak g=q/(1-qy)$ is holomorphic on the disc, since $|qy|<1$ in both
cases.

\emph{Operators.} $q\mapsto\mathcal T(q)$ and $q\mapsto M_0(q)$ are holomorphic on $|q|<1$ as maps
into the bounded operators on $\ell^1$, in operator norm: the row truncations $M_N$ in the proof of
{\cite[Lemma ``compactness and meromorphy'', \textup{lem:fred}]{Bonfioli2026bB}} have finitely many non-zero rows, each a norm-convergent $\ell^\infty$-valued
power series in $q$ (the row $s$ is $e_s\sum_{s'}q^{\max(s,s')}\delta_{s'}$), hence are
holomorphic in norm, and, with $M$ either operator and $\vartheta\in\{0,1\}$ its exponent offset in
\cite[Definition ``section identity'', \textup{def:model}]{Bonfioli2026bB}, $\|M-M_N\|\le\sum_{s>N}2|q|^{\vartheta+s}\to0$ uniformly on $|q|\le r<1$; a locally
uniform norm limit of holomorphic operator-valued maps is holomorphic. With compactness
({\cite[Lemma ``compactness and meromorphy'', \textup{lem:fred}]{Bonfioli2026bB}}) and invertibility near $q=0$ ({\cite[Lemma ``existence, uniqueness, analyticity'', \textup{lem:modelexist}]{Bonfioli2026bB}}), the analytic
Fredholm theorem makes $(I-\mathcal T)^{-1}$ and $R_0=(I-M_0)^{-1}$ meromorphic on $|q|<1$ in
operator norm, with poles only where $1-\Sigma_1=0$, resp.\ $S_e=0$ ({\cite[Corollary ``singularities of the two gapless resolvents'', \textup{cor:sing}]{Bonfioli2026bB}}).

\emph{Bulk quantities.} $u=(2q^{2b})_b$ is a norm-convergent, hence holomorphic, $\ell^1$-valued
function, so $\Psi=R_0u$ is meromorphic $\ell^1$-valued and $t_1=v^{\!\top}\Psi$ is meromorphic.
$1-\mathfrak gt_1$ tends to $1$ as $x\to0$, so it is not identically zero, and $B_0=1/(1-\mathfrak
gt_1)$, $P=B_0\Psi$, $\mathfrak gS=B_0-1$ and $B_{0z}=1+y(B_0-1)$ are meromorphic
({\cite[Remark ``the regrouping of $\mathcal W_0$'', \textup{M3:rem:W0regroup}]{Bonfioli2026bB}}).

\emph{Boundary vectors and the sectors $k^\ast\ne0$.} Every entry of
\cite[\textup{(eq:app-mu)}--\textup{(eq:app-lambda)}]{Bonfioli2026bB} is a combination of $B_0$, $B_{0z}$ and sums
$\sum_bP_bx^{m_b}$ with $m_b\ge0$, so $|\mu^{(y)}_s|,|E^{(y)}_s|,|\lambda^{(y)}_s|\le
C\,(|B_0|+|B_{0z}|+\|P\|_{\ell^1})$ uniformly in $s$, locally uniformly in $x$ away from the poles
of $B_0$ and $\Psi$. Hence $D_e\lambda^{(y)}$ has $\ell^1$-tails
$\sum_{s>N}|e_s\lambda^{(y)}_s|=O(|q|^{N+2})$, locally uniformly, and is meromorphic
$\ell^1$-valued; so is $(I-\mathcal T)^{-1}D_e\lambda^{(y)}$. Its pairing with the bounded vectors
$\mu^{(y)}$ and $\tfrac12E^{(y)}$ is an absolutely convergent sum, locally uniformly convergent
away from the poles (an $\ell^1$-valued continuous map has uniformly small tails on compact sets),
hence meromorphic. This gives $\mathcal W_\pm$; the prefactor $x^{-1}$ is at most a pole at $x=0$,
and is removed below.

\emph{The sector $k^\ast=0$.} By {\cite[Remark ``the regrouping of $\mathcal W_0$'', \textup{M3:rem:W0regroup}]{Bonfioli2026bB}},
$\mathcal W_0=A_0+B_0A_1+B_0^2A_2$. The $A_i$ are not finite sums: each is an absolutely convergent
double sum, over the left and the right chain kinds $\mathsf E,\mathsf Z,\mathsf N_{\pm,b}$ of
{\cite[\textup{(eq:app-W0)}]{Bonfioli2026bB}} ($b\ge1$), of monomials $x^jy^\chi$ with $j\ge0$, $\chi\in\{0,1\}$, times at most
two entries of $\Psi$. Since $|x^jy^\chi|\le1$, the double sum is dominated by a constant times
$(1+\|\Psi\|_{\ell^1})^2$, converges locally uniformly where $\Psi$ is holomorphic, and is
meromorphic.

\emph{Identification with the $x$-adic series.} Choose $r_0>0$ so small that on $|x|<r_0$ the
operators $\mathcal T$, $M_0$ and $M_0+\mathfrak guv^{\!\top}$ have $\ell^1$-norm below $\tfrac12$
({\cite[Lemma ``existence, uniqueness, analyticity'', \textup{lem:modelexist}]{Bonfioli2026bB}}; the rank-one term has norm at most $|\mathfrak g|\,\|u\|_{\ell^1}$) and
$|\mathfrak gt_1|<\tfrac12$. There every object above is the sum
of a norm-convergent Neumann series ($(I-\mathcal T)^{-1}$, $R_0$, $P=(I-M_0-\mathfrak
guv^{\!\top})^{-1}u$, and $B_0=\sum_n(\mathfrak gt_1)^n$), whose terms are holomorphic on
$|x|<r_0$ and whose Taylor series at $0$ are the terms of the $x$-adic Neumann series of
{\cite[Lemma \textup{M3:lem:conv}]{Bonfioli2026bB}}, of $x$-adic order tending to infinity. If $\sum_nh_n$ converges uniformly
on a disc about $0$, the Taylor coefficients of the sum are the sums of the Taylor coefficients of
the $h_n$ (Weierstrass); when $\operatorname{ord}h_n\to\infty$ each of these sums is finite and is
the $x$-adic sum. Applying this to each Neumann series, to the absolutely convergent pairings and
to the double sums $A_i$, the Taylor series at $0$ of the analytic $\mathcal W(x,y)$ is the
$x$-adic series of {\cite[Theorem ``\textup{(M3$'$)}'', \textup{thm:M3prime}]{Bonfioli2026bB}}, specialised at $y$ (the specialisation
$\Z[y][[x]]\to\C[[x]]$ at a constant $y$, or at $y=x^2$, is $x$-adically continuous and commutes
with the sums). In particular that function is holomorphic at $x=0$. At $y=1$ the series is $V$,
and at $y=x^2$ it is $U$ ({\cite[Corollary \textup{cor:M3prime-growth}]{Bonfioli2026bB}}). A power series that agrees near $0$
with a function meromorphic on the disc is continued by it, uniquely. For $G^T_0$, $\Sigma_0$ and
$1-\Sigma_1\not\equiv0$ are holomorphic on $|q|<1$ ({\cite[Theorem \textup{thm:blocks}]{Bonfioli2026bB}}, first clause).
\end{proof}

\begin{proposition}[the shape of the local form]\label{prop:shape}
Let $q_m$ be a travel pole, let $p\ge1$ be the order of the zero of $1-\Sigma_1$ there
\textup{(}finite, Lemma~\ref{lem:polefinite}\textup{)}, and let $\pi$ be the order of the pole of
$(I-\mathcal T)^{-1}$ there \textup{(}finite, {\cite[Lemma ``compactness and meromorphy'', \textup{lem:fred}]{Bonfioli2026bB}}\textup{)}. Then $\pi\ge p\ge1$
and, for each $y$,
\[
   \mathcal W(x,y)=\frac{\mathcal N_y}{(q-q_m)^{\pi}}+\bigl(\text{poles of order}<\pi\bigr)+H_m ,
   \qquad H_m\ \text{holomorphic at }q_m,
\]
with the leading numerator a perfect square
\begin{equation}\label{eq:junctionpairing}
   \mathcal N_y=\frac{c_\star}{2x_m}\,\Pi_y(q_m)^2,\qquad
   \Pi_y(q_m):=\langle\lambda^{(y)},R\rangle ,\qquad c_\star\ne0,\ x_m=\sqrt{q_m},
\end{equation}
where $R$ spans $\ker(I-\mathcal T(q_m))$, normalised by its sections,
\[
   \sum_{s\ge0}q^{ks}R_s=\Sigma_k(q_m)\quad(k\ge0),\qquad\text{in particular }
   \sum_sR_s=\Sigma_0(q_m),\quad\sum_sq^sR_s=\Sigma_1(q_m)=1 ,
\]
and $L$ spans $\ker(I-\mathcal T(q_m)^{\!\top})$, explicitly $L_s=R_s/e_s$. There are no cross
terms between the sectors. In particular $\Pi_y(q_m)\ne0$ implies that $q_m$ is a pole of $\mathcal W(\,\cdot\,,y)$, of order exactly
$\pi$.
\end{proposition}

\begin{proof}
Write $\mathcal R(q)=(I-\mathcal T(q))^{-1}=\sum_{j=1}^{\pi}A_j(q-q_m)^{-j}+\text{hol}$, with
$A_\pi\ne0$. Multiplying $(I-\mathcal T(q))\mathcal R(q)=I$ by $(q-q_m)^{\pi}$ and letting
$q\to q_m$ gives $(I-\mathcal T(q_m))A_{\pi}=0$, so
$\operatorname{ran}A_{\pi}\subseteq\ker(I-\mathcal T(q_m))$, which is one-dimensional by
{\cite[Lemma ``the null space is at most one-dimensional'', \textup{lem:nullspace}]{Bonfioli2026bB}}; the same step on $\mathcal R(q)(I-\mathcal T(q))=I$ gives
$A_{\pi}(I-\mathcal T(q_m))=0$. Hence $A_\pi x=c_\star\langle L,x\rangle R$ with $c_\star\ne0$,
$R$ spanning the kernel and $L$ the cokernel. The normalisation of $R$ is
{\cite[Lemma ``the null space is at most one-dimensional'', \textup{lem:nullspace}]{Bonfioli2026bB}}: with the sections $G^R_k=\sum_sq^{ks}R_s$, one has $G^R_k=G^R_1\Sigma_k$ and $\Sigma_1(q_m)=1$, so $G^R_1=1$ gives
$G^R_k=\Sigma_k(q_m)$. That $L_s=R_s/e_s$ spans the cokernel follows from
$\mathcal T=D_eK$ with $D_e=\operatorname{diag}(e_s)$ and $K$ symmetric: from $D_eKR=R$ one gets
$KD_e(D_e^{-1}R)=KR=D_e^{-1}R$, that is $\mathcal T^{\!\top}L=L$, and $L\in\ell^\infty$ because
$|R_s|\le|e_s|\,\|R\|_{\ell^1}$; the cokernel is one-dimensional by Fredholm. Substituting the
Laurent expansion into $\mathcal W_++\mathcal W_-$ of \eqref{eq:assembly}, and using that
$\lambda^{(y)},\mathcal W_0$ are holomorphic at $q_m$, the leading coefficient is
$(2x_m)^{-1}\langle\lambda,A_\pi D_e\lambda\rangle=(2x_m)^{-1}c_\star\langle L,D_e\lambda\rangle
\langle\lambda,R\rangle=(2x_m)^{-1}c_\star\langle\lambda,R\rangle^2$, since $D_eL=R$. It is
non-zero exactly when $\Pi_y(q_m)\ne0$, and then $q_m$ is a pole of order $\pi$. Finally
$\langle\mathbf 1,\mathcal RE\rangle=\Sigma_0/(1-\Sigma_1)$ ({\cite[Proposition ``the travel $k$-sum telescopes'', \textup{prop:travelsum}]{Bonfioli2026bB}}) has
a pole of order exactly $p$ at $q_m$, since $\Sigma_0(q_m)\ne0$
(Proposition~\ref{prop:numnonvanish}), while a matrix element of $\mathcal R$ has pole order at
most $\pi$; hence $\pi\ge p$.
\end{proof}

\subsection{The junction pairing \textup{(R-J)}}\label{sec:RJ}

\textup{(R-J)} is the statement that $\Pi_y(q_m)=\langle\lambda^{(y)},R\rangle\ne0$ at infinitely
many travel poles, for $y=1$ and $y=q$. Earlier versions carried it as a hypothesis. It is now a
theorem.

\begin{theorem}[closed form of the junction pairing]\label{thm:RJ}
Let $q_m$ be a travel pole, $x=\sqrt{q_m}$, and $R$ normalised as in
Proposition~\ref{prop:shape}. Then
\begin{align}
   (1-\mathfrak g_Vt_1)\,\Pi_1(q_m)
     &=\frac{2q(1+x)}{1-q}\Bigl(t_1+\frac{1+x}{2x}\Bigr),\label{eq:RJclosed1}\\
   \Pi_q(q_m)
     &=\frac{2q\,B_0(1+x)}{1-x}\Bigl(t_1\frac{1-x+q}{1+q}+\frac1{2x}\Bigr),
     \qquad B_0=\frac1{1-\mathfrak g_Ut_1}. \label{eq:RJclosedq}
\end{align}
Consequently $\Pi_1(q_m)>0$ and $\Pi_q(q_m)>0$ at every travel pole $q_m$.
\end{theorem}

\begin{proof}
Throughout $q=q_m$.

\emph{Reduction to the ungapped bulk vector.} Put $\Psi_0=1$ and let $\Psi=(\Psi_b)_{b\ge1}=R_0u$ be
the solution of the bulk system with $\mathfrak g=0$ and source $u$, that is
\[
   \Psi_b=2q^bF_b,\qquad F_t:=\sum_{b\ge0}q^{\max(b,t)}\Psi_b\qquad(b\ge1),
\]
the term $b=0$ of $F_t$ being the source $q^t$. Then $t_1=v^{\!\top}R_0u=\sum_{b\ge1}q^b\Psi_b$ and,
by {\cite[\textup{(eq:smresolvent)}]{Bonfioli2026bB}}, $P=R_0u/(1-\mathfrak gt_1)=B_0\Psi$, where $B_0=1/(1-\mathfrak gt_1)$
({\cite[Remark ``the bulk sum of the model'', \textup{rem:bulksource}]{Bonfioli2026bB}}). So every bulk weight in \eqref{eq:junctionsym}--\eqref{eq:farjunction}
is $B_0\Psi_b$, with $b=0$ standing for the empty and gap-opening chains. In the variable $q=x^2$,
\[
 \begin{gathered}
   x^{\max(2s,2b)}=q^{\max(s,b)},\qquad x^{\max(2s+2,2b)}=q\,q^{\max(s,b-1)},\\
   x^{\max(2s+1,2b-1)}=x\,q^{\max(s,b-1)},\qquad x^{\max(2s+1,2b+1)}=x\,q^{\max(s,b)},
 \end{gathered}
\]
with $\max(s,-1)=s$. At $y=1$ the slot $B_{0z}$ equals $B_0$, and \eqref{eq:junctionsym},
\eqref{eq:farjunction} combine into
\begin{equation}\label{eq:lambdaPsi}
   \lambda^{(1)}_s=B_0\sum_{b\ge0}\Psi_b\Bigl((1+x)\,q^{\max(s,b)}+(q+x)\,q^{\max(s,b-1)}\Bigr),
\end{equation}
hence $\langle\lambda^{(1)},R\rangle=B_0\bigl((1+x)X+(q+x)Y\bigr)$ with
\[
   X=\sum_{b\ge0}\Psi_b\,(KR)_b,\qquad Y=\sum_{b\ge0}\Psi_b\,G_b,\qquad
   (KR)_b=\sum_{t\ge0}q^{\max(b,t)}R_t,\quad G_b=\sum_{s\ge0}q^{\max(s,b-1)}R_s .
\]
All double sums converge absolutely ($\Psi\in\ell^1$ with $|\Psi_b|=O(q^{2b}|\Psi|_{\ell^1})$,
$R\in\ell^1$), so they may be summed in either order.

\emph{Two facts about $R$.} From $\mathcal TR=R$, $R_s=2q^{s+1}(KR)_s$ for every $s\ge0$; and the
normalisation $\sum_sq^sR_s=\Sigma_1(q_m)=1$ of Proposition~\ref{prop:shape} says $(KR)_0=G_0=1$. In
particular $R_0=2q$.

\emph{The $X$ shift identity.} Summing first over $t$ and splitting off $b=0$,
$X=1+S$ with $S=\sum_{b\ge1}\Psi_b(KR)_b$. Summing first over $b$ instead, $X=\sum_tR_tF_t$. The
term $t=0$ is $R_0F_0=2q(1+t_1)$, since $F_0=\sum_bq^b\Psi_b=1+t_1$; for $t\ge1$,
$R_tF_t=2q^{t+1}(KR)_t\cdot\Psi_t/(2q^t)=q\,\Psi_t(KR)_t$. Hence $X=2q(1+t_1)+qS$, and eliminating
$S$,
\[
   X=1+\frac{X-(1+t_1)R_0}{q},\qquad\text{so}\qquad X=\frac{q(1+2t_1)}{1-q}.
\]

\emph{The $Y$ shift identity.} Splitting off $b=0$ and using $G_b=(KR)_{b-1}=R_{b-1}/(2q^b)$ for
$b\ge1$, $Y=1+\sum_{s\ge0}\Psi_{s+1}R_s/(2q^{s+1})$. Summing first over $b$ instead, and using
$q^{\max(s,b-1)}=q^{-1}q^{\max(s+1,b)}$, $Y=\sum_sR_sF_{s+1}/q=\sum_sR_s\Psi_{s+1}/(2q^{s+2})$. The
two expressions give
\[
   Y=1+qY,\qquad\text{so}\qquad Y=\frac1{1-q}.
\]

\emph{The closed forms.} With $q=x^2$,
$(1+x)X+(q+x)Y=\bigl((1+x)q(1+2t_1)+q+x\bigr)/(1-q)$, and since
$(1+x)q+q+x=x(1+x)^2$ this is \eqref{eq:RJclosed1}. At $y=q$ only the slot $s=0$ of $E^{(q)}$
changes, from $B_0$ to $B_{0z}=1+q\,\mathfrak g_US=B_0\bigl(1-t_1q/(1+q)\bigr)$ (using
$\mathfrak g_US=B_0-1=\mathfrak g_Ut_1B_0$). So
$\Pi_q=B_0\bigl((1+x)X+(q+x)Y\bigr)+(B_{0z}-B_0)R_0$, now with $B_0=1/(1-\mathfrak g_Ut_1)$, and
collecting terms with $(1+x)^2(1-x+q)=(1+x)(1+q)-q(1-q)$ (both equal $1+x+x^3+x^4$) gives
\eqref{eq:RJclosedq}.

\emph{Positivity for $\tau_m\le5\cdot10^{-3}$.} By {\cite[Corollary ``the gate closes'', \textup{app:gate}]{Bonfioli2026bB}}, $|\mathfrak g_Vt_1|\le0.983$,
and $\mathfrak g_U<\mathfrak g_V$. Hence $1-\mathfrak gt_1\ge0.017$ for both $y$, and
$q|t_1|\le0.983(1-q)$, so $|t_1|<0.12$ once $q\ge0.9$ (here $q\ge e^{-0.005}$). Since
$(1+x)/(2x)\ge1$, $1/(2x)>1/2$ and $0<(1-x+q)/(1+q)<1$, both brackets are positive, and so are the
prefactors.

\emph{Positivity at the poles above that threshold.} This part is a computer-assisted certificate:
at $q_1,\dots,q_{12}$, $\Pi_1$ and $\Pi_q$ are enclosed in MPFR interval arithmetic and found positive,
together with $1-\mathfrak gt_1>0$ (\texttt{code/zeta\_probe/rj\_certificates/r35rust/src/bin/r54b.rs}). The pole
index comes from \texttt{code/zeta\_probe/rj\_certificates/r39pole}, an interval-arithmetic zero count
of $1-\Sigma_1$ which isolates exactly $13$ zeros on $[0,0.9988]$ \textup{(}the right endpoint
taken as its binary64 value\textup{)}, each in its own cell. Its tiling of the interval was repaired
after a first review found gaps of order $10^{-77}$ between cells, and the repaired code was rebuilt
and re-run from current source by a second reviewer; the trusted base is MPFR and the tool's own code.
\end{proof}

\noindent Only infinitely many poles are needed for Theorems~\ref{thm:V} and~\ref{thm:U}, and the
analytic part ($\tau_m\le5\cdot10^{-3}$) supplies them; the certificate is needed only for the
statement ``at every travel pole''. The algebraic core is formalised in Lean~4
({\cite[\S\,\textup{sec:RJlean}]{Bonfioli2026bB}}); the eigen-relations of $R$, the gate and the analytic estimates behind it enter
the formal statements as hypotheses and are not formalised.

\begin{corollary}[\textup{(R-J)} at the dominant pole]\label{prop:RJdominant}
$\Pi_1(q_1)>0$ and $\Pi_q(q_1)>0$; this is the case $m=1$ of Theorem~\ref{thm:RJ}.
\end{corollary}

\section{Transcendence of the travel resolvent, of \texorpdfstring{$V$}{V} and of
\texorpdfstring{$U$}{U}}\label{sec:V}

The relaxed series is the travel resolvent dressed by the bulk. The resolvent itself is settled
unconditionally; passing from it to $V$ uses the transfer model (Theorem~\ref{thm:model}) and the
junction pairing (Theorem~\ref{thm:RJ}).

\begin{theorem}[the travel resolvent]\label{thm:travelres}
The travel resolvent $G^T_0=\Sigma_0/(1-\Sigma_1)$ is transcendental over $\Q(x)$. This is
unconditional: it uses no growth hypothesis, no asymptotic theorem, and no statement about the
strand-walk model. With the pole set taken from Proposition~\ref{prop:infpolesalt}, as in the proof below, it
has no numerical input. The second route to the pole set, {\cite[Lemma ``infinitude and accumulation of the travel poles'', \textup{lem:infpoles}]{Bonfioli2026bB}}, consumes only
$|T_2(m\pi)|<1$ from the contour bound of {\cite[Lemma ``absolute bound on $T_2$'', \textup{lem:T2abs}]{Bonfioli2026bB}}, against a computed worst value
$0.027$; that bound is a verified enclosure over the parameter region, not a sampled one.
\end{theorem}

\begin{proof}
$\Sigma_0$ and $1-\Sigma_1$ are holomorphic on $|q|<1$ ({\cite[Theorem \textup{thm:blocks}]{Bonfioli2026bB}}, first clause) and
$1-\Sigma_1\not\equiv0$. The zero set of $1-\Sigma_1$ in $(0,1)$ is infinite and accumulates
exactly at $q=1$ (Proposition~\ref{prop:infpolesalt}\textup{(c)} together with Lemma~\ref{lem:polefinite};
alternatively {\cite[Lemma ``infinitude and accumulation of the travel poles'', \textup{lem:infpoles}]{Bonfioli2026bB}}, which consumes from {\cite[Lemma ``absolute bound on $T_2$'', \textup{lem:T2abs}]{Bonfioli2026bB}} only
$|T_2|<1$). At each such $q_m$ the numerator does not
vanish, by Proposition~\ref{prop:numnonvanish}. Hence every $q_m$ is a pole of $G^T_0$, of order
the vanishing order of $1-\Sigma_1$ there, which is finite. So $G^T_0$ has infinitely many poles in
$|q|<1$ accumulating at $|q|=1$; under $q=x^2$ these are infinitely many poles in $|x|<1$
accumulating at $|x|=1$. An algebraic function over $\Q(x)$ has only finitely many poles in
$|x|<1$, so no algebraic function has such a pole set; hence $G^T_0$ is transcendental. It is not
$D$-finite either, in $q$ or in $x$, and $|q|=1$ is its natural boundary
(Theorems~\ref{thm:nonDfinite} and~\ref{thm:natboundary}): for $D$-finiteness the finiteness of the
pole set of an algebraic function is replaced by Lemma~\ref{lem:odeposes}, which confines the poles
in the disc of a solution of a linear differential equation to the finitely many zeros of its
leading coefficient.
\end{proof}

\begin{theorem}\label{thm:V}
The relaxed growth series $V$ is transcendental over $\Q(x)$.
\textup{(}Standing: Theorems~\ref{thm:model} and~\ref{thm:RJ}.\textup{)}
\end{theorem}

\begin{proof}
By Lemma~\ref{lem:mero}, $V$ is meromorphic on $|x|<1$, and the local forms below are those of
this continuation. By Proposition~\ref{prop:shape}, at a travel pole $q_m$ the relaxed series has the local form
$V=\mathcal N_1(q-q_m)^{-\pi}+(\text{poles of order}<\pi)+H_m$ with $H_m$ holomorphic, $\pi\ge1$,
and $\mathcal N_1=(2x_m)^{-1}c_\star\Pi_1(q_m)^2$ with $c_\star\ne0$. That proposition consumes two
non-vanishing statements, and they are separate: $S_e(q_m)\ne0$, in the quantitative form
$|S_e(q_m)|\ge0.35\sqrt{\tau_m}$ of {\cite[Proposition \textup{prop:selower}]{Bonfioli2026bB}}, which makes the bulk boundary
vectors $\lambda^{(1)},\mu^{(1)}$ holomorphic at $q_m$ and says nothing about vanishing; and
$\Sigma_0(q_m)\ne0$, from Proposition~\ref{prop:numnonvanish}, which pins the pole order
$\pi\ge p\ge1$; and $\Pi_1(q_m)\ne0$ at every travel pole, by Theorem~\ref{thm:RJ}.

So the leading Laurent coefficient is non-zero at every travel pole, and $V$ has a pole at each of
them. The travel poles are infinite in number and accumulate at $q=1$, that is at $|x|=1$
(Proposition~\ref{prop:infpolesalt} and Lemma~\ref{lem:polefinite}); this is the boundary of the disc $|x|<1$ of meromorphy, not of the disc
of convergence $|x|<x_1$, which only the dominant pole bounds. They are locally finite in $(0,1)$
(Lemma~\ref{lem:polefinite}), so an infinite subset of them accumulates at $q=1$ and nowhere else.
Hence $V$ has infinitely many poles in $|q|<1$ accumulating at $|q|=1$; under $q=x^2$ these are
infinitely many poles in $|x|<1$ accumulating at $|x|=1$. An algebraic function over $\Q(x)$ has
only finitely many poles in $|x|<1$, so an infinite pole set accumulating at the boundary is
impossible for one; hence $V$ is transcendental.
\end{proof}

\begin{corollary}\label{cor:beta}
The exponential growth rate of $v_n$ is $\beta_2=1.4916177871\ldots$, with
$\beta_2^2$ the reciprocal of the smallest travel pole $q_1$; the value $\tfrac32$ is excluded.
It needs only the location of the dominant travel pole: not {\cite[Remark ``the sharp size of $T_2$; \textup{\textsc{heuristic}}'', \textup{lem:cos}]{Bonfioli2026bB}}, not
{\cite[Appendix \textup{app:annulus}]{Bonfioli2026bB}}, and not {\cite[Appendix \textup{app:star}]{Bonfioli2026bB}}. It does share with
Theorem~\ref{thm:V} the transfer model \textup{(M)}, which is what places $q_1$ on $V$ rather than
on $G^T_0$ alone, and at the dominant pole the junction pairing is positive
(Corollary~\ref{prop:RJdominant}). The
lower bound $\beta_2$ for $u_n$ and $v_n$ is free even of \textup{(M)}, coming from the
pure-travel subcount \cite{Bonfioli2026a}. The same rate for $u_n$
follows from Proposition~\ref{prop:travelinv} (the former hypothesis \textup{(T)} of
Theorem~\ref{thm:U}, now proved)
together with the positivity of $\Pi_q(q_1)$; $u_n\le v_n$ gives the upper bound
unconditionally.
\end{corollary}

Whether the \emph{number} $\beta_2$ is transcendental is a separate question, open, and with a
clean reduction to a value statement for a Hahn--Exton $q$-Bessel function; it is treated in
{\cite[Appendix \textup{app:beta2}]{Bonfioli2026bB}}, which is independent of everything else in this paper.

\subsection{The true series \texorpdfstring{$U$}{U}}\label{sec:U}

\begin{theorem}\label{thm:U}
The growth series $U$ of $W_{\mathrm{univ}}$ is transcendental over $\Q(x)$.
\textup{(}Standing: Theorems~\ref{thm:model} and~\ref{thm:RJ}.\textup{)}
\end{theorem}

\begin{remark}[the inputs of Theorem~\ref{thm:U}]\label{rem:Uinputs}
The travel poles pass from $V$ to $U$ by Proposition~\ref{prop:travelinv}, whose proof does not use
the metric identity. The factorisation {\cite[\textup{(eq:mobius)}]{Bonfioli2026bB}} with the dictionary {\cite[\textup{(eq:dict)}]{Bonfioli2026bB}} and
the reciprocity $t_0=b_1$ is {\cite[Theorem ``the bulk dictionary at the travel poles'', \textup{thm:L}]{Bonfioli2026bB}}. The residue is the square of $\Pi_q(q_m)$, which
is Theorem~\ref{thm:RJ}. The history of the former inputs \textup{(T)} and \textup{(L)} is
recorded in \cite[\S\S\,\textup{sec:inputT}--\textup{sec:inputL}]{Bonfioli2026bB}.

Independently of \textup{(M)}, \textup{(T)} and \textup{(R-J)}, the half
$|s|<1$ of the gate \eqref{eq:gate} holds at every travel pole. The numerator amplitude estimate
$(\star)$ (at every travel pole
$q_m$, $|P_{12}(q_m)|\,w_m^3<2$) is {\cite[Theorem ``the amplitude estimate $(\star)$'', \textup{thm:star}]{Bonfioli2026bB}} ({\cite[Appendix \textup{app:star}]{Bonfioli2026bB}}):
$|P_{12}|\le0.619\,\tau^{3/2}$ for $\tau\le5\cdot10^{-3}$ and $|P_{12}|w^3\le0.75$ at the six
poles above that threshold. The passage to $|s|<1$ is
{\cite[Corollary ``the gate closes'', \textup{app:gate}]{Bonfioli2026bB}}: $|s|\le0.983$ on the uniform range and $|s|\le0.296$ at the six
exceptional poles, with the denominator lower bound $|S_e|\ge0.63\sqrt\tau$ supplied by
{\cite[Lemma ``the denominator at the pole'', \textup{app:Se}]{Bonfioli2026bB}}. The remaining half $b_0>0$ is proved in the same corollary from a sign chain
through the closed form for $b_0$, which is {\cite[Corollary ``the dictionary'', \textup{cor:dict}]{Bonfioli2026bB}}. The
leading-order picture of {\cite[Remark ``derivation of the $U$-numerator amplitude'', \textup{rem:elemY3}]{Bonfioli2026bB}} (the phase-matched zero giving
$|P_{12}|/\tau^{3/2}\to1/(4\sqrt2)$, a fourfold margin below the threshold $1/\sqrt2$) is
consistent with, and explained by, the second-difference identity of
{\cite[Proposition \textup{prop:P12second}]{Bonfioli2026bB}}.
\end{remark}

\begin{proof}
By Lemma~\ref{lem:mero} at $y=x^2$, $U$ is meromorphic on $|x|<1$, and the local forms below are
those of this continuation. By Proposition~\ref{prop:travelinv} (the former input \textup{(T)}), $U$ inherits the travel poles. By Proposition~\ref{prop:shape} at $y=q$, at each
travel pole $U$ has the local form $\mathcal N_q(q-q_m)^{-\pi}+(\text{lower order})+H_m$ with
$\mathcal N_q=(2x_m)^{-1}c_\star\Pi_q(q_m)^2$ and $c_\star\ne0$; the boundary vectors are holomorphic there
because the gate \eqref{eq:gate} holds at \emph{every} travel pole. Indeed, on the uniform range
$\tau\le5\cdot10^{-3}$, $|P_{12}|\le0.619\,\tau^{3/2}$ ({\cite[Theorem ``the amplitude estimate $(\star)$'', \textup{thm:star}]{Bonfioli2026bB}}) and
$|S_e|\ge0.63\sqrt\tau$ ({\cite[Lemma ``the denominator at the pole'', \textup{app:Se}]{Bonfioli2026bB}}) give
$|s|=\tfrac{q}{1-q}\,|P_{12}|/|S_e|\le\tfrac{\tau}{e^\tau-1}\cdot\tfrac{0.619}{0.63}\le0.983<1$,
and $b_0>0$ follows from the sign chain of the same lemmas through the closed form of
{\cite[Corollary ``the dictionary'', \textup{cor:dict}]{Bonfioli2026bB}}; the six exceptional poles are checked directly
({\cite[Corollary ``the gate closes'', \textup{app:gate}]{Bonfioli2026bB}}). By Theorem~\ref{thm:RJ} the leading coefficient $\mathcal N_q$ is
non-zero at every $q_m$, so $U$ has an uncancelled pole at each of them; these accumulate at
$q=1$ (Proposition~\ref{prop:infpolesalt} and Lemma~\ref{lem:polefinite}), so $U$ has infinitely many poles in
the unit disc accumulating at its boundary and is transcendental over $\Q(x)$, as in
Theorem~\ref{thm:V}. Transcendence does not consume the exceptional-pole enumeration of
{\cite[Theorem ``the amplitude estimate $(\star)$'', \textup{thm:star}]{Bonfioli2026bB}}: the uniform range $\tau\le5\cdot10^{-3}$ alone carries infinitely many
poles (by Proposition~\ref{prop:infpolesalt} and Lemma~\ref{lem:polefinite}, only
finitely many travel poles lie outside it); the enumeration only upgrades
``at infinitely many travel poles'' to ``at every travel pole'' in the gate. For orientation: the
computed values over $m=5,\dots,71$ are $|P_{12}|/\tau^{3/2}\to1/(4\sqrt2)=0.1768$ and
$s\to\tfrac14$, well inside every bound above.
\end{proof}

\begin{remark}[the standing of Theorem~\ref{thm:U}]\label{rem:Ustatus}
The unconditional content of this section is the gate: at every travel pole,
$|P_{12}|\le0.619\,\tau^{3/2}$ ({\cite[Theorem ``the amplitude estimate $(\star)$'', \textup{thm:star}]{Bonfioli2026bB}}), $|S_e|\ge0.63\sqrt\tau$
({\cite[Lemma ``the denominator at the pole'', \textup{app:Se}]{Bonfioli2026bB}}) and hence $|s|\le0.983<1$ ({\cite[Corollary ``the gate closes'', \textup{app:gate}]{Bonfioli2026bB}}), together with the
infinitude and accumulation of the travel poles (Proposition~\ref{prop:infpolesalt},
Lemma~\ref{lem:polefinite}) and the fact that these
facts compose to non-algebraicity ({\cite[Remark ``the assembly is machine-checked'', \textup{rem:uassembly}]{Bonfioli2026bB}}).
The factorisation {\cite[\textup{(eq:mobius)}]{Bonfioli2026bB}}, the reciprocity $t_0=b_1$, the dictionary {\cite[\textup{(eq:dict)}]{Bonfioli2026bB}} and
the half $b_0>0$ of the gate are {\cite[Theorem ``the bulk dictionary at the travel poles'', \textup{thm:L}]{Bonfioli2026bB}}; the junction pairing is Theorem~\ref{thm:RJ}.

\textup{(T)} is now proved. Proposition~\ref{prop:travelinv} follows from
{\cite[Corollary \textup{cor:localzero}]{Bonfioli2026bB}}: a pure-travel element has no gap edge, so by
{\cite[Theorem ``no gap edge forces $c=0$'', \textup{thm:nogap}]{Bonfioli2026bB}} some relaxed-optimal realisation is a single open walk, which bounds
$\ell_T\le\ell_R$; the reverse inequality is immediate. The metric identity
$\ell_T=\ell_R+2c$ is not used.

\textup{(R-J)} is Theorem~\ref{thm:RJ}, and \textup{(M)} is Theorem~\ref{thm:model}. The standing of
Theorem~\ref{thm:U} is therefore that of its inputs: the gate ({\cite[Appendix \textup{app:star}]{Bonfioli2026bB}}, which
contains verified interval-arithmetic enclosures); the metric identity and the faithfulness of the
model (Lean); the assembly \textup{(M3)} (a hand proof, {\cite[Appendix \textup{app:M3prime}]{Bonfioli2026bB}}, with two
independent adversarial reviews and exact series agreement to $x^{26}$); and, for the statement at
every travel pole rather than infinitely many, the interval certificate at the first twelve poles.
The assembly is not formalised.
\end{remark}

\section{Beyond algebraicity: \texorpdfstring{$D$}{D}-finiteness, functional equations and the
natural boundary}\label{sec:beyond}

Theorems~\ref{thm:V} and~\ref{thm:U} exclude algebraic equations. The pole set that proves them
excludes much more, because it is infinite and lies inside the disc $|x|<1$ on which $U$ and $V$ are
meromorphic (Lemma~\ref{lem:mero}), well beyond their radius of convergence $x_1=0.6704\ldots$. This
section records what it excludes. Every statement below about $U$ or $V$ has the standing of
Theorems~\ref{thm:V} and~\ref{thm:U} (Theorems~\ref{thm:model} and~\ref{thm:RJ}); every statement
about the travel resolvent $G^T_0$ alone is unconditional, as Theorem~\ref{thm:travelres} is.

\subsection{Poles of \texorpdfstring{$U$}{U} and \texorpdfstring{$V$}{V} on the negative axis}

\begin{proposition}\label{prop:minusx}
At every travel pole $q_m$, $U$ has a pole at $x=-x_m$ as well as at $x=x_m$, both of order
$\pi$, the order of the pole of $(I-\mathcal T)^{-1}$ at $q_m$. At $x=-x_m$ the junction pairing is
negative:
\[
   \Pi_q\big|_{x=-x_m}=\frac{2q_mB_0(1-x_m)}{1+x_m}\Bigl(t_1\frac{1+x_m+q_m}{1+q_m}
   -\frac1{2x_m}\Bigr)<0,\qquad B_0=\frac1{1-\mathfrak g_Ut_1}.
\]
\textup{(}Standing: Theorems~\ref{thm:model} and~\ref{thm:RJ}; at the six poles with
$\tau_m>5\cdot10^{-3}$, the enumeration standing of {\cite[Theorem ``the amplitude estimate $(\star)$'', \textup{thm:star}]{Bonfioli2026bB}}.\textup{)}
\end{proposition}

\begin{proof}
The travel resolvent $(I-\mathcal T)^{-1}$, the bulk vector $\Psi$, $t_1$ and $B_0$ depend on $x$
only through $q=x^2$; $x$ itself enters $\mathcal W(x,q)$ only through the monomials of
\cite[\textup{(eq:app-mu)}--\textup{(eq:app-E)}]{Bonfioli2026bB} and the prefactor $x^{-1}$. The proofs of
Proposition~\ref{prop:shape} and of \eqref{eq:RJclosedq} use only $q=x^2$, never the sign of $x$
(the reductions $x^{\max(2s+1,2b\pm1)}=x\,q^{\max(\cdot,\cdot)}$ and the final polynomial identities
hold for either root). Near $x=-x_m$ one has $q-q_m=(x+x_m)(x-x_m)$ with the second factor
non-zero, so $U=\mathcal N_q(-2x_m)^{-\pi}(x+x_m)^{-\pi}+(\text{lower order})$ with
$\mathcal N_q=c_\star(-2x_m)^{-1}\Pi_q|_{x=-x_m}^2$, and \eqref{eq:RJclosedq} at $x=-x_m$ is the
displayed formula. The prefactor is positive: $0<x_m<1$, and $1-\mathfrak g_Ut_1>0$ because
$|\mathfrak g_Ut_1|<|\mathfrak g_Vt_1|=|s|<1$ ({\cite[Corollary ``the gate closes'', \textup{app:gate}]{Bonfioli2026bB}}). For the bracket write
$t_1=s(1-q)/q$ and use $(1+x+q)/(1+q)=1+x/(1+x^2)\le\tfrac32$:
\[
   2x_m\,|t_1|\,\frac{1+x_m+q_m}{1+q_m}\;\le\;3|s|\,\frac{1-q_m}{x_m}.
\]
For $\tau_m\le5\cdot10^{-3}$, $|s|\le0.983$ and $(1-q_m)/x_m\le\tau_me^{\tau_m/2}\le0.00502$, so the
right side is below $0.015$. At the six poles with $\tau_m>5\cdot10^{-3}$, $|s|\le0.296$
({\cite[Corollary ``the gate closes'', \textup{app:gate}]{Bonfioli2026bB}}) and $(1-q_m)/x_m\le(1-q_1)/x_1=0.8212$, since $q\mapsto(1-q)/\sqrt q$
decreases; the right side is at most $0.730$. In both cases $|t_1|(1+x_m+q_m)/(1+q_m)<1/(2x_m)$, the
bracket is negative, $\mathcal N_q\ne0$, and $-x_m$ is a pole of order exactly $\pi$.
\end{proof}

\noindent At the dominant pole the bracket is $-0.2155$, and it tends to $-\tfrac12$ along the
ladder. Only the uniform range is needed for infinitely many poles at $-x_m$; the six-pole list
upgrades this to every travel pole.

\begin{proposition}[$V$ at $-x_m$]\label{prop:Vminusx}\label{rem:Vminusx}
At every travel pole $q_m$, $V$ has a pole at $x=-x_m$ of order $\pi$. With
$\beta_m:=t_1-(1-x_m)/(2x_m)$,
\[
   (1-\mathfrak g_Vt_1)\,\Pi_1\big|_{x=-x_m}=\frac{2q_m(1-x_m)}{1-q_m}\,\beta_m ,
\]
and for $\tau=\tau_m\le5\cdot10^{-3}$, with $h=\tau/2$, $\gamma=4\sinh h+2-2\cosh h$ and
$\Lambda=e^{h/2}\sinh h/\bigl(2h(4-\gamma/3)\bigr)$,
\[
   \frac{1-5.14\sqrt\tau}{1+1.20\sqrt\tau}\;\le\;\frac{\beta_m}{\Lambda\tau^2}\;\le\;
   \frac{1+5.14\sqrt\tau}{1-1.20\sqrt\tau},
   \qquad\text{so}\qquad 0.073\,\tau^2\le\beta_m\le0.187\,\tau^2 .
\]
At the six poles with $\tau_m>5\cdot10^{-3}$, direct evaluation gives
$\beta_m=0.1168,\ 1.119\cdot10^{-3},\ 1.362\cdot10^{-4},\ 3.484\cdot10^{-5},\ 1.266\cdot10^{-5},\
5.651\cdot10^{-6}$.
\textup{(}Standing: Theorems~\ref{thm:model} and~\ref{thm:RJ}; at the six poles with
$\tau_m>5\cdot10^{-3}$, the enumeration standing of {\cite[Theorem ``the amplitude estimate $(\star)$'', \textup{thm:star}]{Bonfioli2026bB}}.\textup{)}
\end{proposition}

\noindent Here the gate does not decide the sign, unlike for $U$: $t_1=s(e^\tau-1)$ with
$s\to\tfrac14$, and $(1-x_m)/(2x_m)=(e^{\tau/2}-1)/2=\tfrac\tau4+O(\tau^2)$, so the two terms of
$\beta_m$ cancel at leading order. The proof computes the $O(\tau^2)$ remainder. It uses an exact
second-order $q$-difference equation for the $q$-cosine to eliminate the unknown curvature at
the zero.

\begin{proof}
\emph{The pole.} $V=\mathcal W(x,1)$, and the proof of Proposition~\ref{prop:minusx} applies
verbatim at $y=1$: the local form of Proposition~\ref{prop:shape} and \eqref{eq:RJclosed1} use only
$q=x^2$. Near $x=-x_m$ the leading coefficient is $c_\star(-2x_m)^{-1}\Pi_1|_{x=-x_m}^2$ with
$c_\star\ne0$, and \eqref{eq:RJclosed1} at $x=-x_m$ is the displayed identity. Its prefactor is
positive, and $1-\mathfrak g_Vt_1=1-s>0$ by the gate
({\cite[Corollary ``the gate closes'', \textup{app:gate}]{Bonfioli2026bB}}). So $\Pi_1|_{x=-x_m}\ne0$ as soon as $\beta_m\ne0$, and
then $-x_m$ is a pole of order exactly $\pi$. At the six exceptional poles, $\beta_m>0$ is the
listed evaluation. It remains to bound $\beta_m$ for $\tau\le5\cdot10^{-3}$.

\emph{Notation.} Let $c_j=(-1)^jq^{j^2+j}/(q;q)_{2j}$ and $C(y)=\sum_jc_jy^j$, and write
$\psi(t)=\cos(Ze^t;q^2)=C(Z^2e^{2t})$. Also write $\varphi(t)=e^{-t/2}\psi(t)$,
$w=\sqrt{2/\tau}$, $A=\varphi'(0)=\psi'(0)$ and $M_k=\max_{|t|\le\tau}|\varphi^{(k)}(t)|$. At the
pole, $\psi(0)=0$ (Proposition~\ref{prop:qtrig}), so $\varphi(0)=0$. Moreover $z_0=Ze^{-h}$,
$qZ=Ze^{-2h}$ and $x_m=e^{-h}$.

\emph{Step 1: $\beta_m$ as a symmetric sum.} We have $t_1=P_{12}/S_e$ (\S\ref{sec:notation}),
$S_e=\psi(-h)$ (Proposition~\ref{prop:qtrig}) and $P_{12}=-\tfrac12[\psi(h)+\psi(-h)]$
({\cite[Proposition \textup{prop:P12second}]{Bonfioli2026bB}}). With $(1-x_m)/(2x_m)=(e^h-1)/2$ this gives
\begin{equation}\label{eq:betaE1}
   \beta_m=-\frac{\psi(h)+e^h\psi(-h)}{2\psi(-h)}=-\frac{e^{h/2}\,E_1}{2S_e},\qquad
   E_k:=\varphi(kh)+\varphi(-kh).
\end{equation}

\emph{Step 2: a $q$-difference equation.} The coefficients satisfy
$(1-q^{2j-1})(1-q^{2j})c_j=-q^{2j}c_{j-1}$. Summed against $y^j$, this gives
$C(y)-(1+q^{-1})C(q^2y)+q^{-1}C(q^4y)=-q^2yC(q^2y)$. In the variable $t$, $y\mapsto q^2y$ is
$t\mapsto t-2h$. We shift by $2h$ and use $Z^2=2(e^{2h}-1)$. Then we divide by $e^{t/2+h}$, which
gives, for all real $t$,
\begin{equation}\label{eq:chidiff}
   \varphi(t+2h)+\varphi(t-2h)=\bigl(2\cosh h-4\sinh h\,e^{2t}\bigr)\varphi(t).
\end{equation}
At $t=0$, and after differentiating twice at $t=0$, since $\varphi(0)=0$:
\begin{equation}\label{eq:chiat0}
   E_2=0,\qquad \varphi''(2h)+\varphi''(-2h)=(2-\gamma)\,\varphi''(0)-16\sinh h\,A .
\end{equation}

\emph{Step 3: elimination.} Write $a_n=\varphi^{(n)}(0)/n!$. Taylor's theorem with Lagrange
remainders gives the following:
\begin{itemize}
\item $E_k=2a_2(kh)^2+2a_4(kh)^4+r_k$ with $|r_k|\le2(kh)^6M_6/720$, for $k=1,2$;
\item $\varphi''(2h)+\varphi''(-2h)=4a_2+96a_4h^2+r''$ with $|r''|\le\tfrac43h^4M_6$.
\end{itemize}
The two relations \eqref{eq:chiat0} determine $a_2h^2$ and $a_4h^4$. Substituting them into $E_1$
gives, with $D=4-\gamma/3$, the exact identity
\[
   D\,E_1=4h^2\sinh h\,A+\tfrac14h^2r''-\bigl(\tfrac\gamma{16}+\tfrac D4\bigr)r_2+D\,r_1 .
\]
Since $0<\gamma<12$, we have $0<D<4$, and the remainder is at most
$C_6h^6M_6$ with $C_6=\tfrac\gamma{90}+\tfrac13+4\bigl(\tfrac2{45}+\tfrac1{360}\bigr)$.

\emph{Step 4: the slope.} By \eqref{eq:chiat0}, $\varphi(2h)-\varphi(-2h)=-2\varphi(-2h)
=-2e^h\cos(qZ;q^2)$. At a pole, $\cos(qZ;q^2)=qZ\sin(Z;q^2)$ and
$|\sin(Z;q^2)|\ge\sqrt{1-0.61Z}$ for $Z\le\min(0.2,q/3)$; both are proved in
{\cite[proof of Proposition \textup{prop:selower}]{Bonfioli2026bB}}, and here
$Z\le0.1002$. Since $e^hqZ=z_0$, and Taylor's theorem gives
$|\varphi(2h)-\varphi(-2h)-4hA|\le\tfrac83h^3M_3$, we get
\[
   |A|\ \ge\ \frac{z_0\,|\sin(Z;q^2)|}{\tau}-\frac{\tau^2}6M_3
   \ \ge\ w\Bigl[\sqrt{\tfrac{1-q}\tau}\,\sqrt{1-0.61Z}-\tfrac\tau3\,\frac{M_3}{w^3}\Bigr].
\]

\emph{Step 5: derivative bounds.} By {\cite[Lemma \textup{app:rep2}]{Bonfioli2026bB}} at radius $\rho=z_0e^t$,
$\psi(t)=\operatorname{Re}\mathbb E[F]$. Here $F=(\rho e^{i\theta};q)_\infty^{-1}$,
$\theta=\tfrac\pi2+\sigma g$, $\sigma^2=\tau/2$ and $g$ is standard normal. We have
$\partial_tF=-i\partial_\theta F$. Gaussian integration by parts,
$\mathbb E[f^{(k)}(g)]=\mathbb E[f(g)\,\mathrm{He}_k(g)]$ with $\mathrm{He}_k$ the Hermite
polynomials, then gives
\[
   \psi^{(k)}(t)=\operatorname{Re}\bigl[(-i)^kw^k\,\mathbb E[F\,\mathrm{He}_k(g)]\bigr].
\]
In $\log F=\sum_{n\ge1}(\rho e^{i\theta})^n/(n(1-q^n))$, the real parts are bounded term by term:
\begin{itemize}
\item $n=1$: $-\rho\sin(\sigma g)/(1-q)\le\lambda|g|$, with $\lambda=\rho\sigma/(1-q)$;
\item $n=2$: $-\kappa_2\cos(2\sigma g)\le-\kappa_2+2\kappa_2\sigma^2g^2$, with
$\kappa_2=\rho^2/(2(1-q^2))$, using $\cos2u\ge1-2u^2$;
\item $n\ge3$: at most $\rho^3/\bigl(9q^2(1-q)(1-\rho)\bigr)\le K_\ast:=\rho^3/\bigl(9q^2(1-q)(1-\rho/q)\bigr)$, since
$n(1-q^n)\ge3(1-q^3)\ge9q^2(1-q)$.
\end{itemize}
Hence
$|\psi^{(k)}(t)|\le w^ke^{K_\ast-\kappa_2}\,\mathbb E\bigl[e^{\lambda|g|+\kappa_2\tau g^2}
|\mathrm{He}_k(g)|\bigr]$. From
$\varphi^{(k)}=e^{-t/2}\sum_j\binom kj(-\tfrac12)^{k-j}\psi^{(j)}$ and $1/(2w)=\sigma/2$, the same
bound holds for $\varphi^{(k)}$ with the extra factor $e^{\tau/2}$, and with
\[
   \textstyle\sum_j\binom kj(\sigma/2)^{k-j}\,|\mathrm{He}_j(g)|
\]
in place of $|\mathrm{He}_k(g)|$. For $|t|\le\tau$
and $\tau\le\bar\tau=5\cdot10^{-3}$, the parameters $\lambda,\rho,\sigma,\kappa_2\tau,K_\ast$ are at
most their values at $t=\tau=\bar\tau$. They are increasing in $\tau$ and $t$, and
$\kappa_2\ge e^{-2\bar\tau}/2$. Directed-rounding interval arithmetic evaluates the resulting
Gaussian integrals on $6000$ cells of $g\in[0,14]$. The integrand is bounded beyond $g=14$ by its
value at $14$, because it is decreasing there even after multiplication by $e^g$. The result is
\[
   \max_{|t|\le\tau}|\psi''|\le3.24\,w^2,\qquad M_3\le5.68\,w^3,\qquad M_6\le53.3\,w^6
\]
(\texttt{code/zeta\_probe/nondfinite/vminusx\_certificate.py}, mpmath \texttt{iv} at $50$ digits;
the computed enclosures are $3.2196$, $5.6320$ and $52.193$, rounded up to the stated constants, which are the ones used below).

\emph{Step 6: assembly.} Each quantity below is monotone in $\tau$, so it is evaluated at $\bar\tau$
with the rounded constants. Step 4 gives $|A|\ge0.958\,w$. By Taylor's theorem,
$S_e=\psi(-h)=-hA(1+\delta_2)$ with $|\delta_2|\le hK_2w/(2\cdot0.958)\le1.20\sqrt\tau$,
where $K_2=3.24$. Step 3 and $\sinh h\ge h$ give $E_1=(4h^2\sinh h/D)\,A\,(1+\delta_1)$ with
$|\delta_1|\le C_6\cdot53.3\,\sqrt\tau/(4\sqrt2\cdot0.958)\le5.14\sqrt\tau$. By
\eqref{eq:betaE1},
\[
   \beta_m=\frac{e^{h/2}\,4h^2\sinh h}{2hD}\cdot\frac{1+\delta_1}{1+\delta_2}
   =\Lambda\tau^2\,\frac{1+\delta_1}{1+\delta_2},
\]
which is the two-sided bound. Finally, $\tfrac18<\Lambda\le\Lambda(\bar\tau)\le0.12527$, since
$e^{h/2}$, $\sinh h/h$ and $1/D$ are increasing, and $\sqrt{\bar\tau}=0.0707$. This gives
$0.073\,\tau^2\le\beta_m\le0.187\,\tau^2$. In particular $\beta_m>0$.
\end{proof}

\noindent The six exceptional values and the uniform-range bracket are checked against direct
evaluation at the travel poles in \texttt{code/zeta\_probe/nondfinite/vminusx\_check.py}, in
floating point at $45$ digits. Along
the ladder, $\beta_m/\tau_m^2\to\tfrac18$: the value is $0.12560$ at $m=7$ and $0.12507$ at $m=19$,
and $\beta_m/(\Lambda\tau_m^2)-1$ is below $3\cdot10^{-3}$ there, against the certified $0.49$ at
$\bar\tau$. The error terms are lossy by a factor of about $10^2$; the leading constant is not. The identity
$\cos(qZ;q^2)=qZ\sin(Z;q^2)$ is checked there, and \eqref{eq:chidiff} in
\texttt{vminusx\_diffeq\_check.py}, both to working precision. Only the
uniform range is needed for infinitely many poles of $V$ at $-x_m$.

\subsection{\texorpdfstring{$D$}{D}-finiteness}

\begin{lemma}\label{lem:odeposes}
Let $\Omega\subset\C$ be a domain, let $a_0,\dots,a_r,b$ be holomorphic on $\Omega$ with
$a_r\not\equiv0$, and let $f$ be meromorphic on $\Omega$ with
$a_rf^{(r)}+\cdots+a_1f'+a_0f=b$ off the poles of $f$. Then every pole of $f$ in $\Omega$ is a zero
of $a_r$.
\end{lemma}

\begin{proof}
Let $z_0$ be a pole of $f$ of order $k\ge1$. Then $f^{(j)}$ has a pole of order exactly $k+j$ at
$z_0$, so for $j<r$ the term $a_jf^{(j)}$ has a pole of order at most $k+r-1$, and $b$ is
holomorphic. If $a_r(z_0)\ne0$, the term $a_rf^{(r)}$ has a pole of order exactly $k+r$, which
nothing else in the equation cancels. Hence $a_r(z_0)=0$.
\end{proof}

\begin{theorem}\label{thm:nonDfinite}
$U$ and $V$ are not $D$-finite: neither satisfies a linear differential equation
\[
   \textstyle\sum_{j=0}^ra_j(x)f^{(j)}(x)=b(x)
\]
with polynomial coefficients and $a_r\not\equiv0$, nor even one
whose coefficients $a_j$ and $b$ are holomorphic on a neighbourhood of the closed disc $|x|\le1$.
Equivalently, $(u_n)$ and $(v_n)$ are not P-recursive. The same holds for the travel resolvent
$G^T_0$, as a series in $q$ and as a series in $x$, and for $G^T_0$ it is unconditional.
\textup{(}Standing for $U$ and $V$: Theorems~\ref{thm:model} and~\ref{thm:RJ}.\textup{)}
\end{theorem}

\begin{proof}
Let $f$ be one of the series and suppose it satisfies such an equation. By Lemma~\ref{lem:mero}
(for $G^T_0$, {\cite[Theorem \textup{thm:blocks}]{Bonfioli2026bB}}, first clause) $f$ continues to a meromorphic function on the
open unit disc $D$. The expression $\sum_ja_jf^{(j)}-b$ is then meromorphic on $D$ and vanishes near
$0$, so by the identity theorem it vanishes on $D$ off the poles of $f$. By
Lemma~\ref{lem:odeposes} every pole of $f$ in $D$ is a zero of $a_r$. Since $a_r\not\equiv0$ is
holomorphic on a neighbourhood of the compact closed disc, its zeros there are isolated and hence
finite in number. But $f$ has infinitely many poles in $D$: $V$ at every $x_m$
(Theorem~\ref{thm:V}) and every $-x_m$ (Proposition~\ref{prop:Vminusx}), $U$ at every $\pm x_m$ (Theorem~\ref{thm:U} and
Proposition~\ref{prop:minusx}), and $G^T_0$ at every $q_m$ in the variable $q$ and at every $\pm x_m$
in the variable $x$ (Theorem~\ref{thm:travelres}). This is a contradiction. $D$-finiteness of a
power series is equivalent to P-recursiveness of its coefficients \cite[\S6.4]{StanleyEC2}.
\end{proof}

\noindent The pole argument does not reach differential algebraicity: $1/\cos(\pi/(1-x))$ has
poles at $x=1-2/(2n+1)$, $n\ge1$, infinitely many in $|x|<1$ and accumulating at $x=1$, and satisfies
an algebraic differential equation. Whether $U$ and $V$ are $D$-algebraic is open; we expect not.

\begin{corollary}\label{cor:Wnotdfinite}
The bivariate series $\mathcal W(x,y)\in\Z[y][[x]]\subset\Z[[x,y]]$ is not $D$-finite in $(x,y)$.
\textup{(}Standing: Theorems~\ref{thm:model} and~\ref{thm:RJ}.\textup{)}
\end{corollary}

\begin{proof}
$D$-finiteness is preserved under substitution of algebraic power series without constant term
\cite{Lipshitz1989}. If $\mathcal W$ were $D$-finite, the substitution $y=x^2$ would make
$\mathcal W(x,x^2)=U$ $D$-finite ({\cite[Corollary \textup{cor:M3prime-growth}]{Bonfioli2026bB}}), contrary to
Theorem~\ref{thm:nonDfinite}.
\end{proof}

\subsection{Linear \texorpdfstring{$q$}{q}-difference and Mahler equations}

\begin{proposition}\label{prop:noqdiff}
Let $f$ be $U$ or $V$, as series in $x$, or $G^T_0$, as a series in $q$, and let $a_0,\dots,a_r,b$
be polynomials with $a_0a_r\not\equiv0$.
\begin{enumerate}[label=\textup{(\roman*)},leftmargin=2.2em,itemsep=0.1em,topsep=0.2em]
\item If $Q\in\C^\times$ is not a root of unity, then $\sum_{i=0}^ra_i(x)f(Q^ix)\ne b(x)$.
\item If $k\ge2$ is an integer, then $\sum_{i=0}^ra_i(x)f(x^{k^i})\ne b(x)$.
\end{enumerate}
For $G^T_0$ both statements are unconditional \textup{(}\textup{(ii)} uses the window statement of
\cite[Lemma \textup{lem:infpoles}]{Bonfioli2026bB}, hence the verified bound of
\cite[Lemma \textup{lem:T2abs}]{Bonfioli2026bB}\textup{)}; for $U$ and $V$ the standing is that of
Theorems~\ref{thm:V} and~\ref{thm:U}. When $Q$ is a root of unity the case is open: the pole set can
be invariant under $x\mapsto Qx$, as that of $U$ is under $x\mapsto-x$
\textup{(}Proposition~\ref{prop:minusx}\textup{)}, and the argument below gives nothing.
\end{proposition}

\begin{proof}
Let $D$ be the open unit disc, $P$ the pole set of $f$ in $D$, and $\rho>0$ the radius of
convergence of $f$; $P$ lies in $\rho\le|x|<1$, has finitely many points in each disc $|x|\le R<1$,
and is infinite (proof of Theorem~\ref{thm:nonDfinite}). An equation as in (i) or (ii) holds as an
identity of power series near $0$.

(i) If $|Q|>1$, replacing $x$ by $Q^{-r}x$ and $i$ by $r-i$ gives an equation of the same form with
$Q^{-1}$ in place of $Q$ and new leading and trailing coefficients $a_r(Q^{-r}x)$, $a_0(Q^{-r}x)$; so
we may assume $|Q|\le1$. Then every $f(Q^ix)$ is meromorphic on $D$, and the equation holds on $D$
by the identity theorem. Let $E$ be the finite zero set of $a_0$ in $D$. At $p\in P\setminus E$ the
term $a_0f$ is singular and $b$ is not, so some $f(Q^ix)$ with $1\le i\le r$ is singular at $p$,
that is $Q^ip\in P$. Thus every $p\in P\setminus E$ has a \emph{successor} $Q^ip\in P$, $1\le i\le
r$, and every point of $P$ is the successor of at most $r$ points. Follow successors from any $p\in
P$; along the chain the exponent of $Q$ increases strictly.

If $|Q|<1$ the moduli decrease strictly, so the chain stays in the finite set $P\cap\{|x|\le|p|\}$
and ends, necessarily at a point of $E$. A chain of $n$ steps ends at modulus at most
$|Q|^n|p|<|Q|^n$, and that modulus is at least $\rho$, so $n<N:=\log\rho/\log|Q|$. Hence every
point of $P$ is reached backwards from $E$ in fewer than $N$ steps, and $|P|\le|E|\sum_{n<N}r^n$ is
finite.

If $|Q|=1$ the chain stays on the circle $|x|=|p|$, which meets $P$ in a finite set, and its points
are pairwise distinct because $Q$ is not a root of unity; so it ends in $E$. Hence $P$ lies on the
finitely many circles $|x|=|e|$, $e\in E$, each meeting $P$ in finitely many points, and $P$ is
finite. Both cases contradict the infinitude of $P$.

(ii) Here every $f(x^{k^i})$ is meromorphic on $D$ and the equation holds on $D$. Work on the
positive axis: let $P_+=P\cap(0,1)$ and let $E_+$ be the finite zero set of $a_0$ in $(0,1)$. As in
(i), every $p\in P_+\setminus E_+$ has a successor $p^{k^i}\in P_+$ with $1\le i\le r$; it is
smaller than $p$, so every chain ends in $E_+$, and reading a chain backwards,
$p=e^{k^{-n}}$ for some $e\in E_+$ and some integer $n\ge0$. Hence, for $0<R<1$ and
$\ell=\log(1/R)$,
\[
   \#\bigl(P_+\cap(0,R]\bigr)\;\le\;\sum_{e\in E_+}\#\bigl\{n\ge0:k^n\le\log(1/e)/\ell\bigr\}
   \;\le\;|E_+|\Bigl(1+\log_k\frac L\ell\Bigr),\qquad L=\max_{e\in E_+}\log(1/e),
\]
which grows like $\log(1/\ell)$ as $R\uparrow1$. On the other hand, by {\cite[Lemma ``infinitude and accumulation of the travel poles'', \textup{lem:infpoles}]{Bonfioli2026bB}}
there is a travel pole with $w\in(m\pi,(m+1)\pi)$ for every $m\ge4$, at
$x=e^{-1/w^2}<e^{-1/((m+1)\pi)^2}$, equivalently $q=e^{-2/w^2}<e^{-2/((m+1)\pi)^2}$, and it lies in
$P_+$ (Theorems~\ref{thm:V}, \ref{thm:U}, \ref{thm:travelres}). Distinct windows give distinct
poles, so $P_+\cap(0,R]$ has at least $\lfloor c/(\pi\sqrt\ell)\rfloor-4$ points, with $c=1$ in the
variable $x$ and $c=\sqrt2$ in $q$. This grows like $\ell^{-1/2}$, faster than $\log(1/\ell)$, a
contradiction for $R$ close to $1$.
\end{proof}

\noindent The count in (ii) is made on the real axis on purpose. It uses only the travel poles that
are proved to exist, and it needs no statement about where on the circle the poles of $f$
accumulate; compare Remark~\ref{rem:polesdensity}. Mahler equations are those of
\cite{Nishioka1996}. The catalytic equation of \S\ref{sec:setup} is a $q$-difference equation in
the catalytic variable, not in $x$, and is not affected by (i).

\subsection{The natural boundary}

We use the following rationality theorem.

\begin{theorem}[P\'olya--Bertrandias]\label{thm:PB}
Let $\Omega\subset\C$ be a simply connected domain containing $0$ and stable under complex
conjugation, and let $\varphi$ be the Riemann map of the unit disc onto $\Omega$ with
$\varphi(0)=0$, $\varphi'(0)>0$. If $\varphi'(0)>1$ and $f\in\Z[[x]]$ is the Taylor series at $0$ of
a function meromorphic on $\Omega$, then $f$ is a rational function.
\end{theorem}

\noindent This is P\'olya's theorem \cite{Polya1928}, which replaced the disc of Borel's
rationality theorem by a domain measured by transfinite diameter; Bertrandias
\cite{Bertrandias1964} extended it to $p$-adic and adelic domains, whence the name. Bost and
Chambert-Loir prove a generalisation to curves over number fields
\cite[Thm.~7.9]{BostChambertLoir2009}, from which they recover the P\'olya--Bertrandias criterion
\cite[Example~7.10]{BostChambertLoir2009}, and Cantor's rationality theorem
\cite[Thm.~5.2.2]{Cantor1980} is a capacity-theoretic generalisation. Meromorphy suffices, but the strict
inequality $\varphi'(0)>1$ is needed: for $\Omega$ the unit disc, $\varphi'(0)=1$, and
$\sum_nx^{2^n}$ is not rational.

\begin{theorem}\label{thm:natboundary}
\begin{enumerate}[label=\textup{(\alph*)},leftmargin=2.2em,itemsep=0.1em,topsep=0.2em]
\item The travel resolvent $G^T_0=\Sigma_0/(1-\Sigma_1)\in\Z[[q]]$ has $|q|=1$ as a natural
boundary. This is unconditional.
\item $U$ and $V$ have $|x|=1$ as a natural boundary. \textup{(}Standing: Theorems~\ref{thm:model}
and~\ref{thm:RJ}.\textup{)}
\end{enumerate}
More precisely, in each case there is no point $\zeta$ of the unit circle and no disc $B$ about
$\zeta$ such that the function continues meromorphically to $D\cup B$, $D$ the open unit disc.
\end{theorem}

\begin{proof}
Let $f$ be $G^T_0$ (in $q$), $U$ or $V$ (in $x$). In each case $f$ has integer coefficients (the
travel-run counts, $u_n$, $v_n$), is meromorphic on $D$ (Lemma~\ref{lem:mero}), and is not rational,
having infinitely many poles in $D$ (Theorems~\ref{thm:travelres}, \ref{thm:V}, \ref{thm:U}).
Suppose $f$ continues meromorphically to $D\cup B$ with $B$ the disc of radius $\delta>0$ about
$\zeta$, $|\zeta|=1$. Shrinking $\delta$, we may assume $\delta<|\operatorname{Im}\zeta|$ when
$\zeta\notin\R$, so that $B$ and its conjugate $\bar B$ are disjoint; when $\zeta=\pm1$,
$\bar B=B$. Because the coefficients are real, $\overline{f(\bar x)}=f(x)$ near $0$, hence on $D$,
and $x\mapsto\overline{f(\bar x)}$ continues $f$ meromorphically to $\bar B$, agreeing with $f$ on
$\bar B\cap D$. So $f$ is meromorphic and single-valued on
\[
   \Omega=D\cup B\cup\bar B ,
\]
which is stable under conjugation and simply connected: it is the disc $D$ with one or two disjoint
discs attached, each meeting $D$ in a convex set. Let $\varphi$ be its Riemann map, $\varphi(0)=0$,
$\varphi'(0)>0$. By uniqueness of the Riemann map, $\overline{\varphi(\bar z)}=\varphi(z)$, so
$\varphi$ has real Taylor coefficients (Schwarz reflection). The map $\psi=\varphi^{-1}|_D$ sends $D$
into $D$ with $\psi(0)=0$ and is not onto, since $\Omega\ne D$; by the Schwarz lemma
$0<\psi'(0)<1$, that is $\varphi'(0)>1$. Theorem~\ref{thm:PB} now makes $f$ rational, a
contradiction.
\end{proof}

\noindent A $D$-finite series continues analytically along every path avoiding the finitely many
zeros of the leading coefficient of its equation, so Theorem~\ref{thm:natboundary} gives a second
proof of the polynomial-coefficient case of Theorem~\ref{thm:nonDfinite}. The first proof is
elementary and also covers coefficients holomorphic near the closed disc.

\begin{remark}[where on the circle the poles accumulate]\label{rem:polesdensity}
Theorem~\ref{thm:natboundary} does not locate the singularities on the circle. The poles proved
above accumulate only at $x=\pm1$ (at $q=1$ for $G^T_0$), yet every point of the circle is singular.
Numerically the pole set is much larger. The function $1-\Sigma_1=\cos(Z;q^2)$ has $64$ zeros in
$|q|<0.94$ by the argument principle: the travel poles $q_1,q_2$ and $62$ non-real zeros in $31$
conjugate pairs, the first at $q=-0.4075\pm0.5449i$. At each non-real zero $\Pi_q$ was found
non-zero at $40$ digits, so these are candidate poles of $G^T_0$ and of $U$ off the real axis. Plotted,
the zeros approach roots of unity along straight lines. We conjecture that the poles are dense on
the circle; a route is a $q$-cosine resummation at $q=\zeta e^{-t}$, $\zeta$ a root of unity. The
natural boundary does not depend on this.
\end{remark}

\section{Open questions}\label{sec:conclusion}

The travel resolvent $\Sigma_0/(1-\Sigma_1)$ is transcendental, not $D$-finite, and has $|q|=1$ as a
natural boundary, unconditionally. The growth series $U$ of $W_{\mathrm{univ}}$ and its relaxed
companion $V$ have the same properties in $x$, on the transfer model (Theorem~\ref{thm:model}) and
the junction pairing (Theorem~\ref{thm:RJ}); two finite ranges are computer-assisted, the six gate
poles above $\tau=5\cdot10^{-3}$ and the positivity of the pairing at the first twelve poles, and
neither is needed for transcendence, only for the statements ``at every travel pole''. The series
in question is the growth series of the universal group $W_{\mathrm{univ}}$; for a particular right
triangle it agrees with that of $W_T$ only to the effective universality depth of
\cite{Bonfioli2026a}. Several questions remain open.
\begin{enumerate}[label=\textup{(\arabic*)},leftmargin=2.2em,itemsep=0.1em,topsep=0.2em]
\item Are the poles of $U$ dense on the unit circle (Remark~\ref{rem:polesdensity})?
\item Are $U$ and $V$ $D$-algebraic? The pole argument does not decide this.
\item Is the number $\beta_2$ transcendental? The question reduces to a value statement for a
Hahn--Exton $q$-Bessel function, and the obstruction sits at the arithmetic ratio $\rho=\tfrac12$
\cite[Appendix \textup{app:beta2}]{Bonfioli2026bB}.
\item Is $(u_n\bmod p)$ non-automatic for some prime $p$? By Christol's theorem this would give a
second, analysis-free proof of transcendence; the $p$-kernel is computed to be maximally
non-automatic at $p=3,5$ \cite[\S\,\textup{sec:modp}]{Bonfioli2026bB}.
\item The number of travel poles in a window of $w$, equivalently an asymptotic count of the zeros
of $\cos(\cdot\,;q^2)$ in $(0,Z(q))$ uniform as $q\to1$, is not known
\cite[Remark ``what the known theory does not give'', \textup{rem:knowntheory}]{Bonfioli2026bB}.
\end{enumerate}


\begin{thebibliography}{99}
\bibitem{BellMishna2020} J.~P.~Bell and M.~Mishna, \emph{On the complexity of the cogrowth
sequence}, J.~Comb.\ Algebra \textbf{4} (2020), 73--85.

\bibitem{Bertrandias1964} F.~Bertrandias, \emph{Diam\`etre transfini dans un corps valu\'e.
Application au prolongement analytique}, S\'eminaire Delange--Pisot--Poitou, Th\'eorie des
nombres \textbf{5} (1963--1964), expos\'e no.~3, 13~pp.

\bibitem{Bonfioli2026a} V.~Bonfioli, \emph{The universal right-triangle reflection sequence:
metric, universality, and the lamplighter structure}, preprint (2026),
\texttt{doi:10.5281/zenodo.20370090}.

\bibitem{Bonfioli2026b} V.~Bonfioli, \emph{Transcendence of a planar reflection-group growth
series and its relaxed companion}, archival full version (2026), \texttt{doi:10.5281/zenodo.20370090}.
The present paper and its companion are a split of this version; theorem labels are shared with it.

\bibitem{Bonfioli2026bB} V.~Bonfioli, \emph{The growth series of the generic right-triangle
reflection group: block assembly, analytic estimates and certificates}, companion preprint (2026),
\texttt{doi:10.5281/zenodo.20370090}.

\bibitem{Bonfioli2026c} V.~Bonfioli, \emph{Universality and a dimensional transcendence dichotomy
for orthoscheme reflection groups}, preprint (2026), \texttt{doi:10.5281/zenodo.20370090}.

\bibitem{BostChambertLoir2009} J.-B.~Bost and A.~Chambert-Loir, \emph{Analytic curves in algebraic
varieties over number fields}, in: Algebra, Arithmetic, and Geometry (in honor of Yu.~I.~Manin),
Vol.~I, Progr.\ Math.~\textbf{269}, Birkh\"auser, 2009, 69--124; Thm.~7.9 and Example~7.10.

\bibitem{Cannon1984} J.~W.~Cannon, \emph{The combinatorial structure of cocompact discrete
hyperbolic groups}, Geom.\ Dedicata \textbf{16} (1984), 123--148.

\bibitem{Cantor1980} D.~G.~Cantor, \emph{On an extension of the definition of transfinite
diameter and some applications}, J.~Reine Angew.\ Math.\ \textbf{316} (1980), 160--207;
Thm.~5.2.2.

\bibitem{Carlson1921} F.~Carlson, \emph{\"Uber Potenzreihen mit ganzzahligen Koeffizienten},
Math. Z. \textbf{9} (1921), 1--13.

\bibitem{DuchinShapiro2019} M.~Duchin and M.~Shapiro, \emph{The Heisenberg group is pan-rational},
Adv.\ Math.\ \textbf{346} (2019), 219--263.

\bibitem{Exton1978} H.~Exton, \emph{A basic analogue of the Bessel--Clifford equation},
J\~n\=an\=abha \textbf{8} (1978), 49--56.

\bibitem{Grigorchuk1984} R.~I.~Grigorchuk, \emph{Degrees of growth of finitely generated groups
and the theory of invariant means}, Izv.\ Akad.\ Nauk SSSR Ser.\ Mat.\ \textbf{48} (1984),
939--985.

\bibitem{Hahn1953} W.~Hahn, \emph{Die mechanische Deutung einer geometrischen
Differenzengleichung}, Z. Angew.\ Math.\ Mech.\ \textbf{33} (1953), 270--272.

\bibitem{Humphreys1990} J.~E.~Humphreys, \emph{Reflection Groups and Coxeter Groups}, Cambridge
Studies in Advanced Mathematics~29, Cambridge Univ.\ Press, 1990; \S5.12 (Poincar\'e series of a
Coxeter group).

\bibitem{IsolaPiergallini} S.~Isola and R.~Piergallini, \emph{On the generic triangle group},
arXiv:1405.1881 [math.MG], 2014.

\bibitem{KoelinkSwarttouw1994} H.~T.~Koelink and R.~F.~Swarttouw, \emph{On the zeros of the
Hahn--Exton $q$-Bessel function and associated $q$-Lommel polynomials}, J. Math. Anal. Appl.
\textbf{186} (1994), 690--710.

\bibitem{KS1992} T.~H.~Koornwinder and R.~F.~Swarttouw, \emph{On $q$-analogues of the Fourier and
Hankel transforms}, Trans.\ Amer.\ Math.\ Soc.\ \textbf{333} (1992), 445--461.

\bibitem{Lipshitz1989} L.~Lipshitz, \emph{$D$-finite power series}, J.~Algebra \textbf{122} (1989),
353--373.

\bibitem{Nishioka1996} Ku.~Nishioka, \emph{Mahler Functions and Transcendence}, Lecture Notes in
Mathematics~1631, Springer, 1996.

\bibitem{Niven1956} I.~Niven, \emph{Irrational Numbers}, Carus Mathematical Monographs~11,
Mathematical Association of America, 1956; Cor.~3.12 (the rational values of $\cos$ at rational
multiples of $\pi$).

\bibitem{Daalhuis1994} A.~B.~Olde Daalhuis, \emph{Asymptotic expansions for $q$-gamma, $q$-exponential,
and $q$-Bessel functions}, J. Math. Anal. Appl. \textbf{186} (1994), 896--913; \S4 (the numerically
satisfactory pair and exact $q$-Wronskians for the Hahn--Exton equation).

\bibitem{Parry1992} W.~Parry, \emph{Growth series of some wreath products}, Trans.\ Amer.\ Math.\
Soc.\ \textbf{331} (1992), 751--759.

\bibitem{Polya1928} G.~P\'olya, \emph{\"Uber gewisse notwendige Determinantenkriterien f\"ur die
Fortsetzbarkeit einer Potenzreihe}, Math.\ Ann.\ \textbf{99} (1928), 687--706.

\bibitem{Simon} B.~Simon, \emph{Trace Ideals and Their Applications}, 2nd ed., Math.\ Surveys
Monogr.~120, Amer.\ Math.\ Soc., Providence, RI, 2005.

\bibitem{OEIS} N.~J.~A.~Sloane et al., \emph{The On-Line Encyclopedia of Integer Sequences},
\url{https://oeis.org}, sequence A396406.

\bibitem{StampachStovicek2014} F.~\v{S}tampach and P.~\v{S}\v{t}ov\'{\i}\v{c}ek, \emph{The
Hahn--Exton $q$-Bessel function as the characteristic function of a Jacobi matrix}, Spec.\
Matrices \textbf{2} (2014), 131--147.

\bibitem{StanleyEC2} R.~P.~Stanley, \emph{Enumerative Combinatorics, Volume~2}, Cambridge Univ.
Press, 1999.

\bibitem{Stoll1996} M.~Stoll, \emph{Rational and transcendental growth series for the higher
Heisenberg groups}, Invent.\ Math.\ \textbf{126} (1996), 85--109.

\bibitem{Swarttouw1992} R.~F.~Swarttouw, \emph{The Hahn--Exton $q$-Bessel function}, PhD thesis,
Delft University of Technology, 1992; eq.~(4.3.8) (the closed-form $q$-Wronskian).

\end{thebibliography}

\end{document}
